\subsection{\texorpdfstring{映射 \((u, v) \mapsto v \circ u\) 的亚连续性}{映射 (u, v) \textbackslash mapsto v \textbackslash circ u 的亚连续性}}

命题 9. --- 设 R、S、T 是三个 Hausdorff 局部凸空间。假设空间 \(\mathcal{L}(\mathbf{R};\mathbf{S})\) 、\(\mathcal{L}(\mathbf{S};\mathbf{T})\) 、\(\mathcal{L}(\mathbf{R};\mathbf{T})\) 分别赋予简单（相应地，紧、有界）收敛拓扑。则双线性映射 \((u, v) \mapsto v \circ u\) （从 \(\mathcal{L}(\mathrm{R}; \mathrm{S}) \times \mathcal{L}(\mathrm{S}; \mathrm{T})\) 到 \(\mathcal{L}(\mathrm{R}; \mathrm{T})\) 中）是 \((\mathfrak{S}, \mathfrak{T})\) -亚连续的，其中 \(\mathfrak{T}\) 是 \(\mathcal{L}(\mathrm{S}; \mathrm{T})\) 的等度连续子集所成的族，而 \(\mathfrak{S}\) 是 \(\mathcal{L}(\mathrm{R}; \mathrm{S})\) 的有限（相应地，紧、有界）子集所成的族。

我们先证明 \((u, v) \mapsto v \circ u\) 是 \(\mathfrak{T}\) -亚连续的。设 H 是 \(\mathcal{L}(S; T)\) 中的一个等度连续集，W 是 T 中 0 的一个邻域，M 是 R 的一个有限（相应地，紧、有界）子集。我们必须证明，存在 S 中 0 的一个邻域 V，使得若 \(u(M) \subset V\) 且 \(v \in H\) ，则 \(v(u(M)) \subset W\) 。但为此，只需 \(v(V) \subset W\) 对一切 \(v \in H\) 成立，而这样的邻域的存在性由 H 的等度连续性得出。

为了看出 \((u, v) \mapsto v \circ u\) 是 \(\mathfrak{S}\) -亚连续的，我们将证明：对 0 的每个邻域 \(W\) （在 \(T\) 中）、每个有限（相应地，紧、有界）子集 \(M\) （在 \(R\) 中）以及每个有限（相应地，紧、有界）子集 \(L\) （在 \(\mathcal{L}(R; S)\) 中），存在一个有限（相应地，紧、有界）子集 \(N\) （在 \(S\) 中） ，使得关系式 \(v(N) \subset W\) 与 \(u \in L\) 蕴含 \(v(u(M)) \subset W\) 。显然，只需证明可取 \(N = \bigcup_{u \in L} u(M)\) ，

即证明只要 L 与 M 是有限（相应地，紧、有界）的，集 N 就是有限（相应地，紧、有界）的。若 L 与 M 都是有限的，这立即可得；若 M 在 R 中有界且 L 在 \(\mathcal{L}(\mathrm{R};\mathrm{S})\) 中有界（对有界收敛拓扑，cf.~III, 第 22 页），这也立即可得。最后，我们证明：若 M 是 R 中的紧集，且 L 对紧收敛拓扑是 \(\mathcal{L}(\mathrm{R};\mathrm{S})\) 中的紧集，则 N 是 S 中的紧集。但若 \(u_{M}\) 是 \(u \in L\) 到 M 上的限制，则映射 \(u \mapsto u_{M}\) （从 L 到由 M 到 S 中的全体连续映射所成的空间 \(\mathcal{C}(\mathrm{M};\mathrm{S})\)（赋一致收敛拓扑）中）是连续的；因此 L 在此映射下的像是紧的，而我们的断言由映射 \((w,x) \mapsto w(x)\) （从 \(\mathcal{C}(\mathrm{M};\mathrm{S}) \times \mathrm{M}\) 到 S 中）的连续性得出（GT, X, § 1, No.~6, 命题 9）。

在下面两个推论中，我们同命题 9 一样假设：空间 \(\mathcal{L}(\mathbf{R};\mathbf{S})\) 、\(\mathcal{L}(\mathbf{S};\mathbf{T})\) 、\(\mathcal{L}(\mathbf{R};\mathbf{T})\) 三者都赋予简单收敛拓扑，或三者都赋予紧收敛拓扑，或三者都赋予有界收敛拓扑。

推论 1. --- 对 \(\mathcal{L}(\mathbf{S};\mathbf{T})\) 的每个等度连续子集 H，映射 \((u,v)\mapsto v\circ u\) （从 \(\mathcal{L}(\mathbf{R};\mathbf{S})\times \mathbf{H}\) 到 \(\mathcal{L}(\mathbf{R};\mathbf{T})\) 中）是连续的。

这由命题 9（III, 第 32 页）与命题 4（III, 第 31 页）立即得出。

推论 2. --- 假设 S 是桶型空间。若序列 \((u_n)\) 收敛于 \(u\) （在 \(\mathcal{L}(\mathbf{R};\mathbf{S})\) 中），且序列 \((v_{n})\) 收敛于 \(v\) （在 \(\mathcal{L}(\mathbf{S},\mathbf{T})\) 中），则序列 \((v_{n}\circ u_{n})\) 收敛于 \(v\circ u\) （在 \(\mathcal{L}(\mathbf{R};\mathbf{T})\) 中） 。

事实上，序列 \((v_{n})\) 在 \(\mathcal{L}(\mathrm{S};\mathrm{T})\) 中简单有界，由于 \(\mathrm{S}\) 是桶型空间（III, 第 25 页, 定理 1），它是等度连续的；于是该推论是推论 1 的结论。

\section{BOREL 图像定理}

\subsection{Borel 图像定理}

定理 1. --- 设 E 是一个作为 Banach 空间的归纳极限的局部凸空间，F 是一个 Souslin 局部凸空间，例如一个 Lusin 空间（GT, IX, § 6, No.~2 与 No.~4），u 是从 E 到 F 中的一个线性映射。若 u 的图像是 E × F 的一个 Borel 子集，则 u 连续。

设 \(E_i\) 是一族 Banach 空间，\((u_i)\) 是一族连续线性映射 \(u_i: E_i \to E\) ，使得 \(E\) 的拓扑是使诸 \(u_i\) 连续的最细局部凸拓扑。只需证明诸复合映射 \(u \circ u_i\) 都是连续的，或者事实上（GT, IX, § 2, No.~6, 命题 10）证明 \(u \circ u_i\) 到每个满足第一可数公理的闭子空间 \(G\) （在 \(E_i\) 中）上的限制都是连续的。这一限制的图像是 \(u\) 的图像在连续映射 \(u_i \times \text{Id}_F: G \times F \to E \times F\) 下的原像，因而是 \(G \times F\) 中的一个 Borel 集。此外，\(G \times F\) 是一个 Souslin 空间，而 Souslin 空间的每个 Borel 子集都是 Souslin 空间（GT, IX, § 6, No.~3, 命题 10）。于是定理 1 由 GT, IX, § 6, No.~8 的定理 4 得出。

注记. --- 回忆（III, 第 12 页），每个同调的 Hausdorff 半完备空间，例如每个 Fréchet 空间，都是 Banach 空间的归纳极限。* 这对自反 Fréchet 空间的强对偶也同样成立（IV, 第 23 页, 命题 4）。*

\subsection{局部凸 Lusin 空间}

命题 1. --- 设 E 是一个 Hausdorff 局部凸空间。假设存在一个由满足第一可数公理的 Fréchet 空间构成的序列 \((\mathrm{E}_{n})_{n\in\mathbb{N}}\) ，以及连续线性映射 \(u_{n}:E_{n}\to E\) ，使得 \(\mathrm{E}=\bigcup_{n\in\mathbb{N}}u_{n}(\mathrm{E}_{n})\) 。则 E 是一个 Lusin 空间。

设 \(P_n\) 是 \(u_n\) 的核；则 \(u_n\) 定义从商空间 \(E_n / P_n\) 到 \(u_n(E_n)\) 上的一个连续双射。由于 \(E_n / P_n\) 是满足第一可数公理的 Fréchet 空间（GT, IX, § 3, No.~1），从而是一个波兰空间（GT, IX, § 6, No.~1, 定义 1），故 \(u_n(E_n)\) 是 E 的一个 Lusin 子空间（GT, IX, § 6, No.~4, 命题 11）。于是由 GT, IX, § 6, No.~7, 定理 3 的推论，作为正则空间（GT, III, § 3, No.~1）的 E 是一个 Lusin 空间。

例 1. --- 每个满足第一可数公理的 Fréchet 空间都是一个波兰空间，从而是一个 Lusin 空间。因此，空间 \(\mathcal{C}(\mathrm{X})\) （其中 \(\mathrm{X}\) 局部紧且具有可数基）亦然（\(\mathcal{C}(\mathrm{X})\) 的拓扑取为紧收敛拓扑，cf.~GT, X, § 3, No.~3, 推论 与 § 1, No.~6, 推论 3）；* 空间 \(\mathcal{C}^{\infty}(\mathrm{U})\) （其中 U 是 \(\mathbf{R}^{n}\) 的开子集（III, 第 9 页））与空间 \(\mathcal{H}(\mathrm{U})\) （其中 U 是 \(\mathbf{C}^{n}\) 的开子集（III, 第 10 页））也都如此。

命题 1 表明，空间 \(\mathcal{C}_0^\infty (\mathbf{U})\) （其中 \(\mathbf{U}\) 是 \(\mathbf{R}^n\) 中的开集）、\(\mathcal{G}_s(\mathbf{I})\) （其中 \(\mathbf{I}\) 是 \(\mathbf{R}\) 中的紧区间且 \(s\geqslant 1\) ）以及 \(\mathcal{H}(\mathbf{K})\) （其中 \(\mathbf{K}\) 是 \(\mathbf{C}^n\) 的紧子集）都是 Lusin 空间（III, 第 10 页）。

定理 2. --- 设 E 是一个局部凸空间，它是 E 的子空间所成的递增序列 \((\mathrm{E}_n)_{n\in \mathbb{N}}\) 的归纳极限，诸子空间赋予满足第一可数公理的 Fréchet 空间拓扑。假设 E 的每个紧子集都含于某个 \(E_{n}\) 之中，且在这个空间中是紧的。设 F 是满足第一可数公理的 Fréchet 空间。则空间 \(\mathcal{L}_{c}(E;F)\) 是一个 Lusin 空间。

空间 E 是有界型空间（III, 第 12 页），因此空间 \(\mathcal{L}_c(\mathrm{E};\mathrm{F})\) 是完备的（III, 第 23 页, 命题 12）。线性映射 \(j\colon f\mapsto (f|\mathrm{E}_n)_{n\in \mathbf{N}}\) 是从 \(\mathcal{L}_c(\mathrm{E};\mathrm{F})\) 到乘积空间 \(\prod_{n\in \mathbf{N}}\mathcal{L}_c(\mathrm{E}_n;\mathrm{F})\) 中的一个单射；凭借关于 E 的紧子集的假设以及诸 \(\mathfrak{S}\) -拓扑的定义，\(j\) 是从 \(\mathcal{L}_c(\mathrm{E};\mathrm{F})\) 到它的像上（该像赋予由乘积拓扑诱导的拓扑）的一个同构；此外，由于 \(\mathcal{L}_c(\mathrm{E};\mathrm{F})\) 完备，这个像是 \(\prod_{n\in \mathbf{N}}\mathcal{L}_c(\mathrm{E}_n;\mathrm{F})\) 的一个闭子空间

（GT, II, § 3, No.~4, 命题 8）。于是由 GT, IX, § 6, No.~4，只需证明诸空间 \(\mathcal{L}_c(\mathrm{E};\mathrm{F})\) 中的每一个都是 Lusin 空间。在证明的其余部分，我们将假设 E 是满足第一可数公理的 Fréchet 空间。

由于 F 是满足第一可数公理的 Fréchet 空间，它同构于 Banach 空间的一个可数乘积的闭子空间，其中每个 Banach 空间 \(F_{n}\) 都是 F 的商（II, 第 5 页），因而满足第一可数公理。线性映射 \(j': f \mapsto (\mathrm{pr}_{n} \circ f)_{n \in \mathbb{N}}\) 是从 \(\mathcal{L}_{c}(\mathrm{E}; \mathrm{F})\) 到乘积空间 \(\prod_{n \in \mathbb{N}} \mathcal{L}_{c}(\mathrm{E}; \mathrm{F}_{n})\) 中的一个单射，并且利用诸 S-拓扑以及乘积中开集的定义，\(j'\) 是从 \(\mathcal{L}_{c}(\mathrm{E}; \mathrm{F})\) 到它的像上的一个同构；此外，由于 \(\mathcal{L}_{c}(\mathrm{E}; \mathrm{F})\) 完备，这个像是 \(\prod_{n \in \mathbb{N}} \mathcal{L}_{c}(\mathrm{E}; \mathrm{F})\) 的一个闭子空间。因此只需证明诸空间 \(\mathcal{L}_{c}(\mathrm{E}; \mathrm{F}_{n})\) 中的每一个都是 Lusin 空间（GT, IX, § 6, No.~4），从而我们可以假设 F 是满足第一可数公理的 Banach 空间。

空间 \(\mathcal{L}_{c}(\mathrm{E};\mathrm{F})\) 是等度连续闭子集所成可数族的并（III, 第 19 页, 推论 1 与 GT, X, § 2, No.~3, 命题 6）。但 \(\mathcal{L}_{c}(\mathrm{E};\mathrm{F})\) 的每个等度连续子集 H 都是可度量化的且满足第一可数公理（III, 第 18 页, 命题 6 与 GT, X, § 2, No.~4, 定理 1）；若 H 是闭的，则它对由 \(\mathcal{L}_{c}(\mathrm{E};\mathrm{F})\) 的一致结构诱导的一致结构而言是完备的空间，因为后者是完备的。换言之，H 是一个波兰空间，从而更是 Lusin 空间；于是正则空间 \(\mathcal{L}_{c}(\mathrm{E};\mathrm{F})\) 是一个 Lusin 空间（GT, IX, § 6, No.~7, 定理 3 的推论）。

推论. --- 在关于 E 的假设与定理 2 相同的前提下，再假设 E 的每个有界子集都是相对紧的。则 E 的强对偶是一个 Lusin 空间。* 特别地，满足第一可数公理的 Fréchet 空间若同时是 Montel 空间，则其强对偶是 Lusin 空间。*

* 例 2. --- 设 U 是 \(\mathbf{R}^{n}\) 的一个开子集。上述推论特别适用于 Fréchet 空间 \(E = \mathcal{C}^{\infty}(\mathrm{U})\) ；它的对偶 \(\mathcal{C}_{0}^{-\infty}(\mathrm{U})\) （U 上具有紧支集的广义函数所成的空间）于是是一个 Lusin 空间。空间 \(\mathcal{C}_{0}^{\infty}(\mathrm{U})\) 是满足第一可数公理的一列 Fréchet 空间 \(\mathcal{C}_{K_n}^{\infty}(\mathrm{U})\) 的严格归纳极限（III, 第 9 页）。我们可以证明，诸空间 \(\mathcal{C}_{K_n}^{\infty}(\mathrm{U})\) 中的每一个都是 Montel 空间；此外，\(\mathcal{C}_{0}^{\infty}(\mathrm{U})\) 的每个有界子集都含于诸空间 \(\mathcal{C}_{K_n}^{\infty}(\mathrm{U})\) 之一（III, 第 5 页, 命题 6）。于是可以应用定理 2 的推论。从而 \(\mathcal{C}_{0}^{-\infty}(\mathrm{U})\) （\(\mathcal{C}_{0}^{\infty}(\mathrm{U})\) 的对偶，U 上的广义函数空间）对强拓扑而言是一个 Lusin 空间。*

类似地可以证明：对 \(C^{n}\) 的每个开子集 U 以及 \(C^{n}\) 的每个紧子集 K，\(\mathcal{H}(U)\) 的强对偶与 \(\mathcal{H}(K)\) 的强对偶都是 Lusin 空间。*

注记. --- 设 E 如定理 2 中所示；设 F 是一个 Hausdorff 局部凸空间，它是一列连续线性映射 \(u_n: \mathrm{F}_n \to \mathrm{F}\) 的像的并，其中每个 \(\mathrm{F}_n\) 都是满足第一可数公理的 Fréchet 空间；则 \(\mathcal{L}_c(\mathrm{E};\mathrm{F})\) 是一个 Lusin 空间。如命题 1 中那样，我们先化归到每个 \(u_n\) 都是单射的情形；然后，如定理 2 的证明中那样，可以假设 E 是满足第一可数公理的 Fréchet 空间。于是由 I, 第 20 页, 命题 1，\(\mathcal{L}(\mathrm{E};\mathrm{F})\) 是诸 \(\mathcal{L}(\mathrm{E};\mathrm{F}_n)\) 的并；此外，典范单射 \(\mathcal{L}_c(\mathrm{E};\mathrm{F}_n) \to \mathcal{L}_c(\mathrm{E};\mathrm{F})\) 连续（GT, X, § 1, No.~4, 命题 3）。由于由定理 2，诸空间 \(\mathcal{L}_c(\mathrm{E};\mathrm{F}_n)\) 中的每一个都是 Lusin 空间，故 \(\mathcal{L}(\mathrm{E};\mathrm{F}_n)\) 对由 \(\mathcal{L}_c(\mathrm{E};\mathrm{F})\) 的拓扑诱导的拓扑而言也是 Lusin 空间（GT, IX, § 6, No.~4, 命题 11）；于是由 GT, IX, § 6, No.~7, 定理 3 的推论，\(\mathcal{L}_c(\mathrm{E};\mathrm{F})\) 是一个 Lusin 空间。

\subsection{\texorpdfstring{* Banach 空间上的可测线性映射 \(^{1}\)}{Banach 空间上的可测线性映射 \^{}\{1\}}}

命题 2. --- 设 E 是一个 Banach 空间，F 是一个局部凸空间，u 是从 E 到 F 中的一个线性映射。假设对 F 的每个闭子集 B、E 的每个紧子集 X 以及 X 上的每个测度 \(\mu\) ，交集 \(X \cap u^{-1}(B)\) 都是 \(\mu\) -可测的。则 u 连续。

先假设 F 是基域。对 E 的每个紧子集 X 以及 X 上的每个测度 \(\mu\) ，u 到 X 上的限制是 \(\mu\) -可测的（INT, IV）。假设 u 不连续。则可以找到 E 中的点列 \((x_{n})\) ，使得 \(\sum_{n} \|x_{n}\| < \infty\) 且对每个整数 n 有 \(|u(x_{n})| \geqslant n\) 。考虑映射 \(g:(t_{n}) \mapsto \sum_{n} t_{n} x_{n}\) （从立方体 \(C = [0, 1]^{N}\) 到 E 中）；显然 g 连续。因此 \(f = u \circ g\) 对 C 上的每个测度都可测（INT, V）；特别地，对作为 C 的诸因子上 Lebesgue 测度的乘积的测度 \(\mu\) 可测。于是存在 C 的一个紧子集 D，使得 \(\mu(D) > \frac{1}{2}\) ，且 f 在 D 上的限制连续，从而也有界。设 M 是 \(|f|\) 在 D 上的上确界，设 \(p \in N\) 使得 \(p \geqslant 4M\) 。设 \(s = (s_{n})\) 与 \(t = (t_{n})\) 是 D 中满足 \(s_{n} = t_{n}\) （对一切 \(n \neq p\) ）的两个点。则

\[
f (s) - f (t) = u (\sum_ {n} s _ {n} x _ {n} - \sum_ {n} t _ {n} x _ {n}) = (s _ {p} - t _ {p}) u (x _ {p}).
\]

由于 \(|f(s) - f(t)| \leqslant 2\mathbf{M}\) 且 \(|u(x_p)| \geqslant p \leqslant 4\mathbf{M}\) ，我们得到

\[
\left| s _ {p} - t _ {p} \right| \leqslant \frac {1}{2}.
\]

Lebesgue-Fubini 定理（INT, V, 第 2 版, § 8, No.~3, 命题 7 的推论 2）蕴含 \(\mu(D) \leqslant \frac{1}{2}\) ；这就导出矛盾。因此 \(u\) 连续。

在一般情形，对每个 \(v \in F'\) ，线性型 \(v \circ u\) 由前面的论证是连续的。设 \((x_{n})_{n \in \mathbb{N}}\) 是 E 中收敛于 0 的点列；则序列 \((u(x_{n}))_{n \in \mathbb{N}}\) 在 F 中收敛于 0（当 F 赋予拓扑 \(\sigma(\mathrm{F}, \mathrm{F}')\) 时）；因此这个序列对 \(\sigma(\mathrm{F}, \mathrm{F}')\) 有界，从而在 F 中有界（III, 第 27 页, 推论 3）。由于 E 是有界型空间（III, 第 12 页, 命题 2），线性映射 \(u: E \to F\) 连续。*

1 本节的结果依赖于积分（Integration）一书。

\exercisesheading
\exercisegroup

\begin{enumerate}
\def\labelenumi{\arabic{enumi})}
\tightlist
\item
  设 \(\mathbf{E}\) 是非离散拓扑域 \(\mathbf{K}\) 上的一个左拓扑向量空间。称 \(\mathbf{B}\) （\(\mathbf{E}\) 的子集）是有界的，如果对 0 的每个邻域 \(\mathbf{V}\) （\(\mathbf{E}\) 中），存在 \(\lambda \neq 0\) （\(\mathbf{K}\) 中），使得 \(\lambda \mathbf{B} \subset \mathbf{V}\) 。
\end{enumerate}

\begin{enumerate}
\def\labelenumi{\alph{enumi})}
\item
  证明：若 B 在 E 中有界，则对 E 中 0 的每个邻域 V，存在 K 中 0 的一个邻域 U，使得 U.B ⊂ V。
\item
  证明：E 中有界集的闭包是有界的。两个有界集的并是有界的。E 中每个预紧集都是有界的。把 III, 第 4 页, 命题 4 的诸推论推广到 K 上的拓扑向量空间。
\item
  证明：若 A 是 K 中的有界集（把 K 视为自身上的左向量空间），B 是 E 中的有界集，则 A.B 在 E 中有界。
\item
  把 III, 第 4 页 的命题 3 推广到 K 是可度量化拓扑除环的情形。e) 把拟完备空间的概念及其性质推广到拓扑向量空间。
\end{enumerate}

\begin{enumerate}
\def\labelenumi{\arabic{enumi})}
\setcounter{enumi}{1}
\tightlist
\item
  \begin{enumerate}
  \def\labelenumii{\alph{enumii})}
  \tightlist
  \item
    设 E 是非离散拓扑域 K 上的一个左拓扑向量空间。证明：若存在 E 中 0 的一个有界邻域 V（习题 1），则诸集 \(\lambda\) V（\(\lambda \in \mathbf{K}\) 且 \(\lambda \neq 0\) ）构成 E 中 0 的一个基本邻域系。若 K 可度量化，则与 E 的拓扑相伴的 Hausdorff 拓扑（GT, III, § 2, No.~6）可度量化。若 \(\mathbf{K} = \mathbf{R}\) 或 \(\mathbf{K} = \mathbf{C}\) ，则 E 上比 E 的给定拓扑更粗的诸拓扑之中的最细局部凸拓扑（II, 第 80 页, 习题 23）可由单个半范数定义。
  \end{enumerate}
\end{enumerate}

\begin{enumerate}
\def\labelenumi{\alph{enumi})}
\setcounter{enumi}{1}
\item
  证明：局部凸 Hausdorff 空间（其中没有一个仅是点 0）的无穷乘积的拓扑不能由单个半范数定义。
\item
  设 E 是一个局部凸空间，其拓扑由半范数的递增序列 \((p_n)\) 定义。要使 E 的拓扑可由单个半范数定义，必须且只需存在整数 \(n_0\) ，使得对每个 \(n \geqslant n_0\) ，存在数 \(k_n \geqslant 0\) ，使得 \(p_n(x) \leqslant k_n p_{n_0}(x)\) 对一切 \(x \in E\) 成立。
\end{enumerate}

\(d\) ) 设 E 是 \(\mathbf{R}\) 上的向量空间，由区间 \(\mathrm{I} = [0,1]\) 上无穷可微的数值函数组成。对每个整数 \(n\geqslant 0\) ，令 \(p_n(f) = \sup_{0\leqslant k\leqslant n}\left(\sup_{x\in I}|f^{(k)}(x)|\right)\) （其中

\(f^{(0)} = f)\) ）；证明诸 \(p_n\) 是 E 上的范数，且由范数序列 \(p_n\) 定义的拓扑不能由单个范数定义。

\begin{enumerate}
\def\labelenumi{\arabic{enumi})}
\setcounter{enumi}{2}
\item
  设 E 是 R 上的可度量化向量空间，d 是与 E 的拓扑相容的平移不变距离。令 \(|x| = d(x, 0)\) （I, 第 16 页）。证明：对每个整数 n \textgreater{} 0，有 \(\frac{1}{n}|x| \leqslant \left|\frac{x}{n}\right|\) 。由此推出：若 B 是 E 的有界子集，则 \(\sup_{x \in B} |x| < +\infty\) （换言之，B 对距离 d 有界（GT, IX, § 2, No.~3））。给出一个可度量化向量空间 E 以及 E 中一个集合的例子，该集合无界但对距离 d 有界（习题 2）。
\item
  设 E 是非离散可度量化拓扑域 K 上的一个拓扑向量空间。证明：若 E 是 Baire 空间，且 E 中由有界子集构成的有界结构具有可数基（III, 第 37 页, 习题 1），则存在 E 中 0 的一个有界邻域，从而与 E 的拓扑相伴的 Hausdorff 拓扑可度量化（III, 第 37 页, 习题 2）（与习题 6 比较）。
\item
  设 \(\mathbf{E}\) 是非离散赋值域 \(\mathbf{K}\) 上的可度量化向量空间。证明：若 \((\mathbf{B}_n)\) 是 \(\mathbf{E}\) 的有界子集的任意一个序列（III, 第 37 页, 习题 1），则存在标量的一个序列 \((\lambda_n)\) （诸标量 \(\neq 0\)），使得诸集 \(\lambda_n\mathbf{B}_n\) 的并是有界的。
\item
  设 \(\mathrm{E}\) 是一个局部凸空间，它是局部凸 Hausdorff 空间的严格递增序列 \((\mathrm{E}_n)\) 的严格归纳极限，且每个 \(\mathrm{E}_n\) 在 \(\mathrm{E}_{n+1}\) 中闭（II, 第 32 页, 命题 9）。
\end{enumerate}

\begin{enumerate}
\def\labelenumi{\alph{enumi})}
\tightlist
\item
  证明 E 不可度量化（利用 III, 第 5 页, 命题 6 以及前面的习题 5）。\\
\item
  要使 E 的典范有界结构具有可数基，必须且只需每个 \(\mathrm{E}_n\) 的典范有界结构在 \(\mathrm{E}_n\) 中具有可数基。
\end{enumerate}

\begin{enumerate}
\def\labelenumi{\arabic{enumi})}
\setcounter{enumi}{6}
\tightlist
\item
  \begin{enumerate}
  \def\labelenumii{\alph{enumii})}
  \tightlist
  \item
    设 E 是一个无穷维 Banach 空间，S 是 E 的紧凸均衡子集所成的族，它对关系 \(\subset\) 是有向集。证明：E 是 Banach 空间 \(E_{A}\) （A 取遍 S）的归纳系统的归纳极限。（用反证法证明：对诸 \(E_{A}\) 的拓扑的归纳极限拓扑而言，0 的邻域 V 含有一个以 0 为中心的球；为此，注意否则将存在 E 中的点列 \((x_{n})\) ，使得 \(\|x_{n}\| \leqslant 1/n^{2}\) ，且不属于 V。）由此推出：E 中存在不含于任何 \(E_{A}\) （\(A \subset S\) ）的有界子集。
  \end{enumerate}
\end{enumerate}

\begin{enumerate}
\def\labelenumi{\alph{enumi})}
\setcounter{enumi}{1}
\tightlist
\item
  设 E 是一个无穷维 Banach 空间。在向量空间 \(\prod_{m=1}^{\infty}E_{m}\) （其中对每个 m 有 \(E_{m}=E\) ）上，用 \(T_{n}\) 表示这样的拓扑：在每个 \(E_{m}\) （\(m\leqslant n\) ）上取 Banach 空间拓扑的乘积，而对 m\textgreater n 取 \(E_{m}\) 上的最细局部凸拓扑；用 \(F_{n}\) 表示局部凸空间 \(\prod_{m=1}^{\infty}E_{m}\) （赋以拓扑 \(T_{n}\) ）。每个恒等映射 \(F_{n}\to F_{n+1}\) 都连续；证明：归纳系统 \(F_{n}\) 的归纳极限空间是空间 F，其中 F 是在 \(\prod_{m=1}^{\infty}E_{m}\) 上赋以各因子的 Banach 空间拓扑之乘积而得到的空间。由此推出：F 中存在在任何 \(F_{n}\) 中都无界的有界子集。
\end{enumerate}

\begin{enumerate}
\def\labelenumi{\arabic{enumi})}
\setcounter{enumi}{7}
\item
  证明：在一个由拓扑向量空间（在 \(\mathbf{R}\) 或 \(\mathbf{C}\) 上，没有一个仅是点 0）作成的无穷乘积的空间中，典范有界结构不具有可数基（首先证明只需对空间 \(\mathbf{R}^{\mathbf{N}}\) 证明这一点，然后利用 III, 第 38 页, 习题 4）。
\item
  设 E 是 R 上由区间 I = \{0, 1\} 上全体调节函数（FVR, I, 第 4 页）组成的向量空间。对每个整数 n \textgreater{} 0，令 \(V_{n}\) 为全体函数 \(f \in E\) （满足 \(\int_{0}^{1} \sqrt{|f(t)|} dt \leqslant 1/n\) ）所成的集。证明：诸集 \(V_{n}\) 构成 E 上一个与 E 的向量空间结构相容的可度量化拓扑的 0 的基本邻域系，且对该拓扑诸集 \(V_{n}\) 是有界的；但每个 \(V_{n}\) 的凸包是整个空间 E。（注意每个函数 \(f \in E\) 都可写成 \(f = \frac{1}{2}(g + h)\) ，其中 g 与 h 属于 E，且
\end{enumerate}

\[
\int_ {0} ^ {1} \sqrt {| g (t) |} d t = \int_ {0} ^ {1} \sqrt {| h (t) |} d t = \frac {1}{\sqrt {2}} \int_ {0} ^ {1} \sqrt {| f (t) |} d t.
\]

由此推出：E 上比 E 的拓扑更粗的唯一的局部凸拓扑是 E 上的最粗拓扑。

\begin{enumerate}
\def\labelenumi{\arabic{enumi})}
\setcounter{enumi}{9}
\tightlist
\item
  设 \((\mathrm{E}_{\iota})_{\iota \in I}\) 是非离散拓扑域 K 上由 Hausdorff 拓扑向量空间组成的一个无穷族，其中没有一个仅是点 0。设 E 是诸 \(\mathbf{E}_{\iota}\) 的直和向量空间，\(\mathcal{T}_0\) 是 I, 第 24 页, 习题 14 中定义的 E 上的拓扑。那么，E 的子集 B 对 \(\mathcal{T}_0\) 有界（III, 第 37 页, 习题 1），当且仅当 B 含于某个乘积子空间 \(\prod_{\iota \in H} \mathbf{E}_{\iota}\) ，其中 H 是 I 的一个有限
\end{enumerate}

子集，且 B 在每个 \(E_{\iota}\) （\(\iota \in H\) ）上的投影是有界的。由此推出（对 K = R 或 C）：若每个 \(E_{\iota}\) 都是拟完备空间，则赋以 \(T_{0}\) 的 E 是拟完备的。

\begin{enumerate}
\def\labelenumi{\arabic{enumi})}
\setcounter{enumi}{10}
\tightlist
\item
  设 \(\mathbf{E}\) 是带非离散赋值的域 \(\mathbf{K}\) 上的一个拓扑向量空间。
\end{enumerate}

\begin{enumerate}
\def\labelenumi{\alph{enumi})}
\item
  为使 E 的均衡子集 A 吸收 E 的每个有界子集（III, 第 37 页, 习题 1），只需 A 吸收 E 中每个趋于 0 的序列 \((x_{n})\) 的点所成的集即可。此时称 A 是吸有界集的。
\item
  设 \(u\) 是从 \(E\) 到一个拓扑向量空间 \(F\) （\(K\) 上的）中的线性映射。\(E\) 的每个有界子集在 \(u\) 下的像在 \(F\) 中有界，当且仅当对每个点列 \((x_{n})\) （\(E\) 中的、趋于 0 的），序列 \((u(x_{n}))\) 在 \(F\) 中有界。
\item
  假设 \(\mathbf{E}\) 可度量化。证明：\(\mathbf{E}\) 的每个吸有界集的子集都是 \(\mathbf{E}\) 中 0 的一个邻域。由此推出：若 \(u\) 是从 \(\mathbf{E}\) 到一个拓扑向量空间 \(\mathbf{F}\) （\(\mathbf{K}\) 上的）中的线性映射，它把 \(\mathbf{E}\) 中每个收敛于 0 的序列变为 \(\mathbf{F}\) 中的有界序列，则 \(u\) 在 \(\mathbf{E}\) 上连续。
\end{enumerate}

\begin{enumerate}
\def\labelenumi{\arabic{enumi})}
\setcounter{enumi}{11}
\item
  设 E 是带非离散赋值的域 K 上的 Hausdorff 拓扑向量空间，F 是 K 上的可度量化向量空间。若 u 是从 E 到 F 中的连续线性映射，使得对 F 的每个有界子集 B，\(u^{-1}(\mathbf{B})\) 在 E 中有界，证明 u 是从 E 到 F 的一个子空间上的同构。
\item
  设 \(I\) 是无穷集，\((E_{\iota})_{\iota\in I}\) 是一族局部凸空间，其中没有一个是 0。设 \(f\) 是从 \(E = \prod_{\iota\in I}E_{\iota}\) 到 Banach 空间 \(F\) 中的一个线性映射。证明：若 E 的每个有界子集在 f 下的像都是 F 的有界子集，则存在 I 的有限子集 H，使得对每个 \(\iota \notin H\) ，f 在 \(E_{\iota}\) （视为 E 的子空间）上的限制是零。（用反证法：否则可以构造 E 中的有界序列 \((x_{n})\) ，使其在 f 下的像在 F 中无界。）
\item
  证明：若可度量化局部凸空间 E 的拓扑不能由单个范数定义，则 E 的典范有界结构不存在可数基（利用 III, 第 38 页, 习题 5，证明否则 E 中将存在一个有界的吸有界集的集合（III, 第 39 页, 习题 11），并利用 III, 第 39 页, 习题 11, c) 完成论证）。
\item
  在 R 上的 Hausdorff 拓扑向量空间 E 中，设 A 是紧凸集，B 是闭的凸有界集。证明：并 \(A \cup B\) 的凸包 C 是闭集。（考虑 C 的闭包中不属于 A 的一点 z，并化归到 z = 0 的情形。注意存在 0 的一个邻域 V 和一个数 \(\alpha < 1\) ，使得关系式 \(0 \leqslant \lambda \leqslant 1\) 、\(x \in A\) 、\(y \in B\) 、\(\lambda x + (1 - \lambda)y \in V\) 蕴含 \(\lambda \leqslant \alpha\) 。然后，对 0 的每个邻域 W，考虑三元组 \((\lambda, x, y)\) （满足 \(\lambda x + (1 - \lambda)y \in W\) 、\(0 \leqslant \lambda \leqslant 1\) 、\(x \in A\) 、\(y \in B\) ）全体所成的集。）
\item
  设 \(\mathbf{E}\) 是满足第一可数公理的可度量化局部凸空间，且其完备化 \(\hat{\mathbf{E}}\) 是满足第一可数公理的 Fréchet 空间。
\end{enumerate}

证明：\(\hat{E}\) 的每个有界子集 B 都含于 \(\hat{E}\) 的某个有界子集的闭包中。（化归到 B 可数、并排成序列 \((x_{n})\) 的情形；另一方面，设 \((p_{n})\) 是定义 \(\hat{E}\) 的拓扑的半范数的递增序列；对每个整数 n，考虑 E 的点列 \((y_{nk})_{k\geqslant1}\) ，它收敛于 \(x_{n}\) ，且满足 \(p_{n}(x_{n}-y_{nk})\leqslant1\) （对一切 \(k\geqslant1\)）。）

\begin{enumerate}
\def\labelenumi{\arabic{enumi})}
\setcounter{enumi}{16}
\tightlist
\item
  设 \(\mathbf{A}\) 是 \(\geqslant 1\) 的、从 \(\mathbf{N}\) 到 \(\mathbf{N}\) 中的递增映射全体所成的集；对每个 \(\alpha \in \mathbf{A}\) ，用 \(\mathbf{B}_{\alpha}\) 表示由所有点 \(z = (z_n) \in \mathbf{R}^{\mathbf{N}}\) （它们满足 \(|z_n| \leqslant \alpha(n)\) ，对一切 \(n \in \mathbf{N}\) ）组成的集。
\end{enumerate}

\begin{enumerate}
\def\labelenumi{\alph{enumi})}
\item
  证明：诸集 \(\mathbf{B}_{\alpha}\) 构成空间 \(\mathbf{R}^{\mathbf{N}}\) 的全体有界子集所成的有界结构的一个基。b) 对每个 \(\alpha \in \mathbf{A}\) ，集 \(\mathbf{RB}_{\alpha}\) 是 \(\mathbf{R}^{\mathbf{N}}\) 的向量子空间，它异于 \(\mathbf{R}^{\mathbf{N}}\) 且在 \(\mathbf{R}^{\mathbf{N}}\) 中稠密；因此存在一个线性型 \(f_{\alpha} \neq 0\) （不连续的，在 \(\mathbf{R}^{\mathbf{N}}\) 上），使得 \(f_{\alpha}(z) = 0\) 对一切 \(z \in \mathbf{B}_{\alpha}\) 成立。
\item
  设 \(\mathbf{E}\) 是由所有映射 \(g: \alpha \mapsto (g_n(\alpha)) \in \mathbf{R}^{\mathbf{N}}\) （从 \(\mathbf{A}\) 到 \(\mathbf{R}^{\mathbf{N}}\) 中，且对一切 \(n \in \mathbf{N}\) 使和 \(p_n(g) = \sum_{\alpha} |g_n(\alpha)|\) 为有限）全体所成的向量空间。证明诸 \(p_n\) 是半范数
\end{enumerate}

它们在 E 上定义出一个 Fréchet 空间的拓扑。

\(d\) ) 设 \(\mathbf{H}\) 是由所有 \(h\in \mathbf{E}\) （满足 \(h(\alpha)\in \mathbf{RB}_2\) ，对一切 \(\alpha \in \mathbf{A}\)）组成的集；证明 \(\mathbf{H}\) 是 \(\mathbf{E}\) 的一个处处稠密的向量子空间（注意：若 \(h\in \mathbf{E}\) 满足 \(h(\alpha) = 0\) （除 \(\alpha \in \mathbf{A}\) 的有限多个值外），则它属于 \(\mathbf{H}\) ）。

\begin{enumerate}
\def\labelenumi{\alph{enumi})}
\setcounter{enumi}{4}
\tightlist
\item
  设 \(\mathrm{E}_0 \subset \mathrm{E}\) 是 \(\mathrm{E}\) 中由所有 \(g \in \mathrm{E}\) （满足 \(\sum_{\alpha} |f_{\alpha}(g(\alpha))| < +\infty\) ）组成的向量子空间；
\end{enumerate}

此时映射 \(u: g \mapsto (f_{\alpha}(g(\alpha)))_{\alpha \in \mathbf{A}}\) 是从 \(\mathrm{E}_0\) 到 Banach 空间 \(\mathrm{F} = \ell^1(\mathrm{A})\) 中的一个线性映射（I, 第 4 页）。证明 \(u(\mathrm{E}_0)\) 在 \(\mathrm{F}\) 中处处稠密（注意对每个有限子集 \(\mathrm{J}\) （\(\mathrm{A}\) 的），存在 \(g \in \mathrm{E}_0\) ，使得 \(g(\alpha) = 0\) （对一切 \(\alpha \in \mathrm{A} - \mathrm{J}\)），且诸 \(f_{\alpha}(g(\alpha))\) （对 \(\alpha \in \mathrm{J}\)）取 \(\mathbf{R}\) 中任意值）。证明 \(u^{-1}(0)\) 在 \(\mathrm{E}_0\) 中处处稠密（利用 \(d\) ）。最后证明：对每个有界子集 \(C\) （\(E\) 的），存在 \(\alpha_0 \in A\) ，使得 \(f_{\alpha_0}(g(\alpha_0)) = 0\) 对一切 \(g \in C \cap E_0\) 成立，并由此推出 \(u(C \cap E_0)\) 在 \(F\) 中的闭包不是 \(F\) 中 0 的邻域。

\begin{enumerate}
\def\labelenumi{\alph{enumi})}
\setcounter{enumi}{5}
\tightlist
\item
  设 G 是 u 在 \(E_{0} \times F\) 中的图像，它是 Fréchet 空间 \(E \times F\) 的一个向量子空间。证明 G 在 \(E \times F\) 中处处稠密（注意对每个 \(x \in E_{0}\) ，\(x + u^{-1}(0)\) 在 E 中稠密）。然而证明：对 G 的每个有界子集 M，闭包 \(\overline{M}\) （M 在 \(E \times F\) 中的）不包含有界集 \(\{0\} \times U\) （\(E \times F\) 中的），其中 U 是 F 中的单位球（若 \(N = \text{pr}_{1}(M)\) ，注意由 e)，\(\overline{u(N)}\) 不可能包含 U）。
\end{enumerate}

\begin{enumerate}
\def\labelenumi{\arabic{enumi})}
\setcounter{enumi}{17}
\tightlist
\item
  在 Banach 空间 \(\ell^1 (\mathbf{N})\) （I, 第 4 页）中，设 \(e_m\) 为这样的序列 \((\delta_{mn})_{n\geq 0}\) ：\(\delta_{mn} = 0\) （对 \(m\neq n\)）且 \(\delta_{nn} = 1\) 。定义一个从 \(\ell^1 (\mathbf{N})\) 到 \(\mathbf{R}\) 中的连续映射，它把由诸 \(e_n\) 组成的有界序列变为 \(\mathbf{R}\) 的一个无界子集（利用 Urysohn 定理（GT, IX, §4, No.~2, 定理 2））。
\end{enumerate}

\exercisegroup

\begin{enumerate}
\def\labelenumi{\arabic{enumi})}
\tightlist
\item
  设 E 是一个局部凸空间，T 是它的拓扑。在 E 上使有界集与 T 的有界集相同的诸局部凸拓扑中，有一个 \(T'\) 比其他所有的都细，且它是这些拓扑中唯一有界型的拓扑。在 E 上赋以 \(T'\) 所得到的空间称为与 E 相伴的有界型空间。从 E 到一个局部凸空间 F 中的线性映射 u 把 E 的每个有界子集变为 F 的有界子集，当且仅当它对拓扑 \(T'\) 连续。
\end{enumerate}

证明：拓扑 \(\mathcal{T}'\) 是 \(E\) 上使诸典范单射 \(E_A \to E\) （其中 \(A\) 取遍 \(E\) 的凸的有界均衡子集的族）都连续的局部凸拓扑中最细的一个。

\begin{enumerate}
\def\labelenumi{\arabic{enumi})}
\setcounter{enumi}{1}
\tightlist
\item
  设 I 是无穷集，\((\mathrm{E}_{\iota})_{\iota \in \mathrm{I}}\) 是一族局部凸空间，其中没有一个是 0。a) 假设每个空间 \(\mathrm{E}_{\iota}\) 都是有界型的。证明：若此外乘积空间 \(\mathbf{R}^{\mathrm{I}}\) 还是有界型的，则乘积 \(\mathrm{E} = \prod_{\iota \in \mathrm{I}}\mathrm{E}_{\iota}\) 是有界型的（利用 III, 第 11 页, 命题 1 (iii)
\end{enumerate}

与第 39 页, 习题 15，化归为证明以下命题：一个线性映射 \(f\) （从 \(E\) 到一个 Banach 空间中），若它把每个有界集变为有界集，且它在每个 \(E_{\iota}\) 上的限制都是零，则它在 \(E\) 上必为零。为此，对每个 \(x = (x_{\iota}) \in E\) ，考虑 \(f\) 在诸直线 \(\mathbf{R}x_{\iota}\) 的乘积上的限制。

\begin{enumerate}
\def\labelenumi{\alph{enumi})}
\setcounter{enumi}{1}
\tightlist
\item
  由 \(a\) ) 推出：有界型空间的序列 \((\mathrm{E}_n)\) 的每个乘积都是有界型的。
\end{enumerate}

\begin{enumerate}
\def\labelenumi{\arabic{enumi})}
\setcounter{enumi}{2}
\tightlist
\item
  设 \(\mathbf{E}\) 是一个局部凸空间，\(\mathbf{L}\) 是 \(\mathbf{E}\) 的一个有限余维数的向量子空间，\(\mathbf{S}\) 是 \(\mathbf{L}\) 中一个凸的均衡的吸有界集的集合（III, 第 39 页, 习题 11）。我们将证明存在一个凸的、均衡的且吸有界集的集合 \(\mathbf{S}'\) （\(\mathbf{E}\) 中），使得 \(\mathbf{S} = \mathbf{S}' \cap \mathbf{L}\) 。
\end{enumerate}

\begin{enumerate}
\def\labelenumi{\alph{enumi})}
\item
  可化归到如下情形：L 是一个超平面，使得 \(E = L \oplus Ra\) （对一点 \(a \notin L\) ），且存在 E 中的有界序列 \((x_{n})\) ，使得若记 \(x_{n} = \lambda_{n}(y_{n} + a)\) ，其中 \(\lambda_{n} \in R\) 、\(y_{n} \in L\) ，则 \(|\lambda_{n}|\) 趋于 \(+\infty\) ；若 \(B_{0}\) 是由 a 与诸 \(x_{n}\) 组成的集的均衡凸包，则对一切 n 有 \(y_{n} + a \in \lambda_{n}^{-1}B_{0}\) 。
\item
  设 \(\mathfrak{B}\) 是 \(E\) 中包含 \(\mathbf{B}_0\) 的有界凸均衡子集全体所成的集；根据假设，对每个 \(\mathbf{B} \in \mathfrak{B}\) ，存在 \(\rho_{\mathbf{B}} > 0\) ，使得 \(2\rho_{\mathbf{B}}\mathbf{B} \cap \mathbf{L} \subset \mathbf{S}\) 。证明：若 \(\mathbf{R}\) 是诸集 \(\rho_{\mathbf{B}}\mathbf{B}\) （\(\mathbf{B} \in \mathfrak{B}\) ）的并的均衡凸包，则有 \(\mathbf{R} \cap \mathbf{L} \subset \mathbf{S}\) 。
\end{enumerate}

\begin{enumerate}
\def\labelenumi{\arabic{enumi})}
\setcounter{enumi}{3}
\tightlist
\item
  由习题 3 推出：若 E 是局部凸的有界型空间，则 E 的每个有限余维数的子空间都是有界型的（cf.~IV, 第 64 页, 习题 11）。
\end{enumerate}

\exercisegroup

\begin{enumerate}
\def\labelenumi{\arabic{enumi})}
\item
  设 X 是一个 Hausdorff 拓扑空间，F 是一个拓扑向量空间（在 R 或 C 上）。证明：在从 X 到 F 中的全体连续映射所组成的空间 \(\mathcal{C}(X; F)\) 上，紧收敛拓扑与向量空间结构相容。
\item
  设 \(\mathbf{E}\) 与 \(\mathbf{F}\) 是两个 Hausdorff 局部凸空间，\(\mathfrak{S}\) 是 \(\mathbf{E}\) 的有界子集的一个族。
\end{enumerate}

\begin{enumerate}
\def\labelenumi{\alph{enumi})}
\item
  证明：若 F 不只是 0，则 \(\mathcal{L}_{\mathfrak{S}}(\mathrm{E};\mathrm{F})\) 为 Hausdorff 的一个必要（且充分）条件是 S 中诸集的并在 E 中是完全的（利用 Hahn-Banach 定理）。
\item
  假设 \(\mathfrak{S}\) 是 E 的一个覆盖。证明：存在从 F 到 \(\mathcal{L}_{\mathfrak{S}}(\mathrm{E};\mathrm{F})\) 的一个闭子空间上的同构。由此推出：若 \(\mathcal{L}_{\mathfrak{S}}(\mathrm{E};\mathrm{F})\) 拟完备，则 F 必拟完备。
\item
  假设 \(\mathfrak{S}\) 是与 E 相适配有界结构（III, 第 3 页, 定义 4）。要使 \(\mathcal{L}_{\mathfrak{S}}(\mathrm{E};\mathrm{F})\) 可度量化，必须且只需 F 可度量化，且有界结构 \(\mathfrak{S}\) 具有可数基（III, 第 1 页）。要使 \(\mathfrak{S}\) -拓扑在 \(\mathcal{L}(\mathrm{E};\mathrm{F})\) 上可由单个范数定义，必须且只需 F 的拓扑可由单个范数定义，且存在一个集 \(\mathbf{M} \in \mathfrak{S}\) ，它吸收 \(\mathfrak{S}\) 中的每个集。
\end{enumerate}

\begin{enumerate}
\def\labelenumi{\arabic{enumi})}
\setcounter{enumi}{2}
\item
  设 E 是 R（相应地，C）上的一个拓扑向量空间。证明：对 R（相应地，C）的、其中没有一个是点 0 的有界子集的每个族 S，空间 \(\mathcal{L}_{\mathfrak{S}}(\mathbf{R};\mathbf{E})\) （相应地，\(\mathcal{L}_{\mathfrak{S}}(\mathbf{C};\mathbf{E})\) ）典范地同构于 E。由此推出：对每个整数 n \textgreater{} 0 以及 \(R^{n}\) （相应地，\(C^{n}\) ）的由有界子集组成的每个覆盖 S，\(\mathcal{L}_{\mathfrak{S}}(\mathbf{R}^{n};\mathbf{E})\) （相应地，\(\mathcal{L}_{\mathfrak{S}}(\mathbf{C}^{n};\mathbf{E})\) ）同构于 \(E^{n}\) 。
\item
  \begin{enumerate}
  \def\labelenumii{\alph{enumii})}
  \tightlist
  \item
    设 \(\mathrm{E}_1, \mathrm{E}_2, \mathrm{F}\) 是三个拓扑向量空间（在 \(\mathbf{R}\) 或 \(\mathbf{C}\) 上）。设 \(f\) 是从 \(\mathrm{E}_1\) 到 \(\mathrm{E}_2\) 中的连续线性映射，\(\mathfrak{S}_1\) （相应地，\(\mathfrak{S}_2\) ）是 \(\mathrm{E}_1\) （相应地，\(\mathrm{E}_2\) ）的有界子集的一个族，使得 \(f(\mathfrak{S}_1) \subset \mathfrak{S}_2\) 。证明 \(u \mapsto u \circ f\) 是从 \(\mathcal{L}_{\mathfrak{S}_2}(\mathrm{E}_2; \mathrm{F})\) 到 \(\mathcal{L}_{\mathfrak{S}_1}(\mathrm{E}_1; \mathrm{F})\) 中的一个连续线性映射。
  \end{enumerate}
\end{enumerate}

\begin{enumerate}
\def\labelenumi{\alph{enumi})}
\setcounter{enumi}{1}
\tightlist
\item
  设 \(\mathbf{E},\mathbf{F}\) 是两个拓扑向量空间，\(\mathbf{M}\) 是 \(\mathbf{E}\) 的一个向量子空间。设 \(f\) 是从 \(\mathbf{E}\) 到 \(\mathrm{E / M}\) 上的典范映射，\(\mathfrak{S}\) 是 \(\mathbf{E}\) 的有界子集的一个族。证明映射 \(u\mapsto u\circ f\) 是从 \(\mathcal{L}_{f(\mathfrak{S})}(\mathrm{E / M};\mathrm{F})\) 到 \(\mathcal{L}_{\mathfrak{S}}(\mathrm{E};\mathrm{F})\) 的如下子空间上的同构：该子空间由从 \(\mathbf{E}\) 到 \(\mathbf{F}\) 中、且在 \(\mathbf{M}\) 上为零的连续线性映射组成。
\end{enumerate}

\begin{enumerate}
\def\labelenumi{\arabic{enumi})}
\setcounter{enumi}{4}
\tightlist
\item
  设 \((\mathrm{E}_{\alpha})_{\alpha\in\mathrm{A}}\) 是一族局部凸空间，E 是一个向量空间（与诸 \(E_{\alpha}\) 有相同的标量基域），且对每个 \(\alpha\in A\) ，设 \(h_{\alpha}\) 是从 \(E_{\alpha}\) 到 E 中的一个线性映射。在 E 上赋以使诸 \(h_{\alpha}\) 连续的最细局部凸拓扑（II, 第 27 页）。对每个 \(\alpha \in \mathrm{A}\) ，设 \(\mathfrak{S}_{\alpha}\) 是 \(\mathrm{E}_{\alpha}\) 的有界子集的一个族，并设 \(\mathfrak{S}\) 是诸族 \(h_{\alpha}(\mathfrak{S}_{\alpha})\) （\(\mathrm{E}\) 的有界子集的族）的并。在这些条件下证明：对每个局部凸空间 \(\mathrm{F}\) ，\(\mathfrak{S}\) -拓扑在 \(\mathcal{L}(\mathrm{E};\mathrm{F})\) 上是使线性映射 \(u\mapsto u\circ h_{\alpha}\) （从 \(\mathcal{L}(\mathrm{E};\mathrm{F})\) 到 \(\mathcal{L}_{\mathfrak{S}_{\alpha}}(\mathrm{E}_{\alpha};\mathrm{F})\) 中）都连续的最粗拓扑。特别地，若 \(\mathrm{E}\) 是族 \((\mathrm{E}_{\alpha})_{\alpha \in \mathrm{A}}\) 的拓扑直和（II, 第 29 页, 定义 2）（每个 \(\mathrm{E}_{\alpha}\) 与 \(\mathrm{E}\) 的一个子空间等同），则乘积空间 \(\prod_{\alpha \in \mathrm{A}}\mathcal{L}_{\mathfrak{S}_{\alpha}}(\mathrm{E}_{\alpha};\mathrm{F})\) 典范地同构于空间 \(\mathcal{L}_{\mathfrak{S}}(\mathrm{E};\mathrm{F})\) ，
\end{enumerate}

其中 \(\mathfrak{S}\) 是诸 \(\mathfrak{S}_{\alpha}\) 在 \(\mathfrak{P}(\mathrm{E})\) 中的并。

\begin{enumerate}
\def\labelenumi{\arabic{enumi})}
\setcounter{enumi}{5}
\item
  设 \((\mathrm{E}_{\iota})_{\iota \in I}\) 是一族 Hausdorff 局部凸空间，其中没有一个是 0，设 E 是乘积空间 \(\prod_{\iota \in I} \mathrm{E}_{\iota}\) ，F 是一个赋范空间。证明：存在从赋以有界收敛拓扑（相应地，简单收敛拓扑、预紧收敛拓扑）的空间 \(\mathcal{L}(\mathrm{E};\mathrm{F})\) 到诸空间 \(\mathcal{L}(\mathrm{E}_{\iota};\mathrm{F})\) 的拓扑直和空间上的一个典范同构，其中每个这样的空间都赋以有界收敛拓扑（相应地，简单收敛拓扑、预紧收敛拓扑）。（注意：若 u 是从 E 到 F 中的一个连续线性映射，则存在 I 的有限子集 H，使得 \(u^{-1}(0)\) 包含诸 \(E_{\iota}\) （对所有指标 \(\iota\notin H\) ）的乘积。）
\item
  设 \(\mathrm{E}\) 、\(\mathrm{F}_1\) 、\(\mathrm{F}_2\) 是三个拓扑向量空间，\(f\) 是从 \(\mathrm{F}_1\) 到 \(\mathrm{F}_2\) 中的一个连续线性映射，\(\mathfrak{S}\) 是 \(\mathrm{E}\) 的有界子集的一个集；证明 \(u \mapsto f \circ u\) 是从 \(\mathcal{L}_{\mathfrak{S}}(\mathrm{E};\mathrm{F}_1)\) 到 \(\mathcal{L}_{\mathfrak{S}}(\mathrm{E};\mathrm{F}_2)\) 中的一个连续线性映射。
\item
  设 \(\mathbf{E}\) 是一个拓扑向量空间，带有由有界子集组成的一个集 \(\mathfrak{S}\) 。设 \((\mathbf{G}_{\iota})_{\iota \in \mathbf{I}}\) 是一族拓扑向量空间，\(\mathbf{F}\) 是一个向量空间（与 \(\mathbf{E}\) 及诸 \(\mathbf{G}_{\iota}\) 有相同的标量基域）；对每个 \(\iota \in \mathbf{I}\) ，设 \(g_{\iota}\) 是从 \(\mathbf{F}\) 到 \(\mathbf{G}_{\iota}\) 中的一个线性映射。假设在 \(\mathbf{F}\) 上赋以使诸 \(g_{\iota}\) 连续的最粗拓扑。 证明：\(\mathfrak{S}\) -拓扑在 \(\mathcal{L}(\mathbf{E};\mathbf{F})\) 上是使线性映射 \(u\mapsto g_{\iota}\circ u\) （从 \(\mathcal{L}(\mathbf{E};\mathbf{F})\) 到 \(\mathcal{L}_{\mathfrak{S}}(\mathbf{E};\mathbf{G}_{\iota})\) 中）都连续的最粗拓扑。 特别地，若 \(\mathrm{F} = \prod_{\iota \in \mathrm{I}}\mathrm{G}_{\iota}\) ，则乘积空间 \(\prod_{\iota \in \mathrm{I}}\mathcal{L}_{\mathfrak{S}}(\mathbf{E};\mathbf{G}_{\iota})\) 典范地
\end{enumerate}

等同于 \(\mathcal{L}_{\mathfrak{S}}(\mathrm{E};\mathrm{F})\)

\begin{enumerate}
\def\labelenumi{\arabic{enumi})}
\setcounter{enumi}{8}
\item
  设 E 与 F 是两个 Hausdorff 拓扑向量空间，H 是 \(\mathcal{L}(\mathrm{E};\mathrm{F})\) 的一个等度连续子集。证明：若 E 中存在一个可数的完全集，且 F 的每个有界子集都可度量化，则 H 对 E 上简单收敛拓扑可度量化。此外，若 F 的每个有界子集都满足第一可数公理，则 H 也满足第一可数公理。
\item
  设 \(\mathbf{E}\) 是一个拓扑向量空间且是 Baire 空间，\(\mathbf{F}\) 是一个拓扑向量空间。
\end{enumerate}

\begin{enumerate}
\def\labelenumi{\alph{enumi})}
\item
  证明：若一个子集 \(\mathrm{H}\) （\(\mathcal{L}(\mathrm{E};\mathrm{F})\) 的）对简单收敛拓扑有界，则 \(\mathrm{H}\) 等度连续（对 0 的每个闭邻域 \(\mathrm{V}\) （\(\mathrm{F}\) 中），考虑诸集 \(\mathbf{M}_n = \bigcap_{u\in \mathrm{H}}u^{-1}(n\mathrm{V})\) ）。
\item
  证明：若 \(\mathcal{L}(\mathrm{E},\mathrm{F})\) 的子集 H 不等度连续，则所有 \(x\in E\) （使 \(\mathrm{H}(x)\) 在 F 中不有界）组成的集是一个第一纲集的余集。由此推出：若 \((\mathrm{H}_{n})\) 是 \(\mathcal{L}(\mathrm{E};\mathrm{F})\) 的不等度连续子集的序列，则存在 \(x\in E\) ，使得诸集 \(\mathrm{H}_{n}(x)\) 中没有一个在 F 中有界（« 奇点凝聚原理 »）。
\end{enumerate}

¶ 11) 设 T 是一个可度量化的拓扑空间，E 是一个拓扑向量空间且是 Baire 空间，M 是从 E × T 到一个拓扑向量空间 F 中的映射的一个族，满足下列条件：

\(1^{\circ}\) 对一切 \(t_0 \in \mathbf{T}\) ，映射 \(x \mapsto f(x, t_0)\) （其中 \(f\) 取遍 \(\mathbf{M}\) ）全体所成的集是从 \(\mathbf{E}\) 到 \(\mathbf{F}\) 中的线性映射的一个等度连续集；

\(2^{\circ}\) 对一切 \(x_0 \in \mathbf{E}\) ，映射 \(t \mapsto f(x_0, t)\) （从 \(\mathbf{T}\) 到 \(\mathbf{F}\) 中，其中 \(f\) 取遍 \(\mathbf{M}\) ）全体所成的集是等度连续的。

证明 \(\mathbf{M}\) 等度连续。（给定 \(t_0 \in \mathbf{T}\) 以及 F 中 0 的一个闭的均衡邻域 \(\mathbf{V}\) ，对每个 \(x \in E\) ，令 \(d_{x}\) 为 T 中以 \(t_{0}\) 为中心且满足下列条件的开球半径的上确界：对这些球中任意一点 t，都有 \(f(x, t) - f(x, t_{0}) \in \mathrm{V}\) （对一切 \(f \in M\)）。证明 \(x \mapsto d_{x}\) 在每一点 \(x_{0} \in E\) 处上半连续；为此，证明：若有 \(d_{x} > \alpha > d_{x_{0}}\) （对任意接近 \(x_{0}\) 的点），那么对 F 中 0 的每个邻域 W，\(f(x_{0}, t) - f(x_{0}, t_{0})\) 都将属于 \(V + W\) （当 \(d(t, t_{0}) \leqslant \alpha\) 且 \(f \in M\) 时）。最后利用 GT, IX, § 5, No.~4, 定理 2。）

\begin{enumerate}
\def\labelenumi{\arabic{enumi})}
\setcounter{enumi}{11}
\tightlist
\item
  设 \(\mathbf{E}\) 是一个有界型局部凸空间，\(\mathfrak{S}\) 是 \(\mathbf{E}\) 的有界子集的一个族，它包含每个收敛于 0 的序列的像。
\end{enumerate}

\begin{enumerate}
\def\labelenumi{\alph{enumi})}
\item
  证明：对每个局部凸空间 \(\mathbf{F}\) ，\(\mathcal{L}_{\mathfrak{S}}(\mathrm{E};\mathrm{F})\) 的每个有界子集都等度连续。
\item
  证明：若 F 是 Hausdorff 的拟完备局部凸空间，则空间 \(\mathcal{L}_{\mathfrak{S}}(E;F)\) 是拟完备的。
\end{enumerate}

\begin{enumerate}
\def\labelenumi{\arabic{enumi})}
\setcounter{enumi}{12}
\tightlist
\item
  证明：若 E 是 Hausdorff 的半完备局部凸空间，则对每个局部凸空间 F，\(\mathcal{L}(\mathrm{E};\mathrm{F})\) 中对简单收敛拓扑有界的每个子集对每个 S-拓扑都有界。
\end{enumerate}

\exercisegroup

\begin{enumerate}
\def\labelenumi{\arabic{enumi})}
\item
  证明：Hausdorff 桶型空间的完备化是桶型的。
\item
  设 E 是 R 或 C 上的一个向量空间。证明：赋以 E 上最细局部凸拓扑（II, 第 25 页）的 E 是桶型的。由此给出不可度量化且不是 Baire 空间的桶型空间的例子。
\item
  设 \(E\) 是一个 Hausdorff 局部凸空间，具有可数无穷基 \((a_{n})\) 。
\end{enumerate}

\begin{enumerate}
\def\labelenumi{\alph{enumi})}
\item
  证明：E 具有一个可数的、拓扑无关的基 \((e_n)\) （利用 E 中每条直线都有拓扑补这一事实，用归纳法定义诸 \(e_n\) ）。
\item
  证明：要使 E 是桶型的，必须且只需 E 的拓扑 T 与 E 上的最细局部凸拓扑相同（注意每个序列 \((\lambda_{n}e_{n})\) 的均衡凸包在 E 中闭）。特别地，若 T 可度量化，则 E 不是桶型的（cf.~习题 2）。
\end{enumerate}

\begin{enumerate}
\def\labelenumi{\arabic{enumi})}
\setcounter{enumi}{3}
\item
  设 E 是一个 Banach 空间，其中存在一个无穷的代数无关序列 \((a_{n})\) ，它在 E 中是完全的（例如空间 \(\ell^{1}(\mathbf{N})\) （I, 第 4 页））。设 B 是 E 的一个包含诸 \(a_{n}\) 的基；我们知道（II, 第 80 页, 习题 24）B 不可数。设 \((e_{n})\) 是 B 中互异且异于诸 \(a_{n}\) 的元素组成的序列，并设 C 是诸 \(e_{n}\) 组成的集在 B 中的余集。设 \(F_{n}\) 是 E 中由 C 与诸 \(e_{k}\) （指标 \(k \leqslant n\) ）生成的向量子空间；E 是诸 \(F_{n}\) 的并。设 S 是 E 中的单位球；证明存在一个指标 n，使得 \(S \cap F_{n}\) 不是第一纲集。由此推出：对这个 n 值，\(F_{n}\) 是一个可度量化的、非完备的 Baire 空间。
\item
  给出一个局部凸空间的例子，它是完备的 Hausdorff Baire 空间，但不可度量化（cf.~GT, IX, § 5, 习题 16）。
\item
  称局部凸空间 \(\mathbf{E}\) 是相对有界的，如果 \(\mathbf{E}\) 中存在一个有界的桶。
\end{enumerate}

\begin{enumerate}
\def\labelenumi{\alph{enumi})}
\item
  要使 E 相对有界，必须且只需 E 的拓扑比某个由半范数定义的拓扑粗。此时存在 E 的典范有界结构的一个由桶组成的基。
\item
  要使 E 既是有界型的又是相对有界的，必须且只需 E 的拓扑是 E 上一族赋范空间拓扑的下确界（cf.~III, 第 40 页, 习题 1）。此外，要使 E 的典范有界结构还具有可数基，必须且只需 E 的拓扑是赋范空间拓扑的一个序列的下确界。
\end{enumerate}

\begin{enumerate}
\def\labelenumi{\arabic{enumi})}
\setcounter{enumi}{6}
\item
  称局部凸空间 E 是次桶型的，如果 E 的每个吸有界集的（III, 第 39 页, 习题 11）桶都是 E 中 0 的邻域。每个有界型空间都是次桶型的；每个桶型空间都是次桶型的。证明：Hausdorff 次桶型空间的完备化是桶型的（利用如下事实：在 Hausdorff 局部凸空间 E 中，每个桶都吸收 E 的每个凸的、均衡的、有界的且半完备的子集）。
\item
  设 \((\mathrm{E}_{\iota})_{\iota\in I}\) 是一族次桶型空间，且对每个 \(\iota\in I\) ，设 \(f_{\iota}\) 是从 \(\mathrm{E}_{\iota}\) 到一个向量空间 E 中的线性映射。证明：赋以使诸 \(f_{\iota}\) 连续的最细局部凸拓扑的空间 E 是次桶型的。特别地，次桶型空间的每个商空间都是次桶型的；次桶型空间的每个拓扑直和都是次桶型的。
\item
  设 I 是一个不可数无穷集；在直和向量空间 \(E = R^{(I)}\) 上，一方面考虑最细局部凸拓扑 T，另一方面考虑 I, 第 24 页, 习题 14 中定义的拓扑 \(T_{0}\) ，它是局部凸的；我们知道 T 与 \(T_{0}\) 不同（II, 第 75 页, 习题 11），且赋以 \(T_{0}\) 的 E 是完备的（GT, III, § 3, 习题 10）。证明：E 中的有界集对 T 与 \(T_{0}\) 是相同的（III, 第 39 页, 习题 10），且赋以 \(T_{0}\) 的 E 不是桶型的（注意由所有 \(x = (\xi_{\mathfrak{u}}) \in \mathbf{E}\) （满足 \(\sum_{\mathfrak{u} \in I} |\xi_{\mathfrak{u}}| \leqslant 1\) ）组成的集 T 是一个桶，并利用 II, 第 75 页, 习题 11）。
\item
  证明：若一个次桶型空间中每个闭的凸均衡有界子集都是半完备的，则它是一个桶型空间。
\item
  设 \(\mathbf{E}\) 是一个次桶型空间，\(\mathbf{F}\) 是一个局部凸空间。证明：\(\mathcal{L}(\mathbf{E};\mathbf{F})\) 中对有界收敛拓扑有界的每个子集都等度连续。
\end{enumerate}

¶ 12) a) 设 E 是一个局部凸空间，\((A_n)\) 是 E 中凸均衡集的递增序列，使得 \(A = \bigcup_{n} A_n\) 是吸收的。设 \((W_n)\) 是凸均衡的

0 的邻域组成的递减序列；那么诸 \(\mathrm{W}_n\cap \mathrm{A}_n\) 的均衡凸包 V 是吸收的；若 E 是桶型的，则 \(\overline{\mathbf{V}}\) 是 0 的一个邻域。

\begin{enumerate}
\def\labelenumi{\alph{enumi})}
\setcounter{enumi}{1}
\item
  设 \(\mathfrak{F}\) 是 E 上的一个滤子；假设对每个 \(n\) ，存在一个集 \(\mathrm{M}_n \in \mathfrak{F}\) ，使得 \((\mathrm{M}_n + \mathrm{W}_n) \cap \mathrm{A}_{2n} = \emptyset\) 。设 \(\mathrm{V}_n\) 是诸 \(\mathrm{W}_k \cap \mathrm{A}_k\) （\(k \leqslant n - 1\) ）与 \(\mathrm{W}_n\) 的均衡凸包，使得 \(\mathrm{V}_n\) 是 0 的一个邻域，且有 \(\frac{1}{2}\overline{\mathrm{V}} \subset \mathrm{V}_n\) （对一切 \(n\)）。证明 \((\mathrm{M}_n + \mathrm{V}_n) \cap \mathrm{A}_n = \emptyset\) （对一切 \(n\) ）。
\item
  由 a) 与 b) 推出：若 E 是桶型的且 \(\mathfrak{F}\) 是 E 上的一个 Cauchy 滤子，则存在整数 N，使得对一切 \(\mathbf{M} \in \mathfrak{F}\) 以及 E 中 0 的每个邻域 W，\(\mathbf{M} + \mathbf{W}\) 都与 \(\mathrm{A}_{\mathbf{N}}\) 相交。（用反证法；采用 b) 的记号，考虑一个集合 \(\mathbf{M} \in \mathfrak{F}\) （具有小阶 \(\frac{1}{2}\overline{\mathrm{V}}\)）。）
\end{enumerate}

¶ 13) a) 设 E 是一个桶型空间，\((C_n)\) 是凸的、均衡的集的递增序列，使得 \(E = \bigcup_{n} C_n\) 。设 U 是一个凸的、均衡的且吸收的集，使得对每个 n，

\(U \cap C_{n}\) 在 \(C_{n}\) 中闭。证明 U 是 E 中 0 的一个邻域。（证明 \(\overline{U} \subset 2U\) ，为此考虑 U 上的一个滤子 \(\mathfrak{F}\) （收敛于一点 \(x \in E\) ）并应用习题 12, c)。）

\begin{enumerate}
\def\labelenumi{\alph{enumi})}
\setcounter{enumi}{1}
\tightlist
\item
  设 \(\mathrm{E}\) 是一个桶型空间，\((\mathrm{E}_n)\) 是 \(\mathbf{E}\) 的子空间的递增序列，使得 \(\mathrm{E} = \bigcup_{n} \mathrm{E}_n\) 。
\end{enumerate}

证明：若 \(\mathbf{U}\) 是 \(\mathbf{E}\) 的一个子集，使得 \(\mathbf{U} \cap \mathbf{E}_n\) 是 \(\mathbf{E}_n\) 中的一个桶（对每个 \(n\) ），则 \(\mathbf{U}\) 是 \(\mathbf{E}\) 中 0 的一个邻域。特别地，\(\mathbf{E}\) 是序列 \((\mathbf{E}_n)\) 的严格归纳极限（II, 第 33 页）。

¶ 14) a) 设 E 是一个 Hausdorff 局部凸空间，L 是 E 的一个有限余维数的子空间，T 是 L 中的一个桶。证明：存在 E 中的一个桶 T'，使得 T' ∩ L = T（证明可以取 T' 为 T 在 E 中的闭包 T̅ 与一个有限维紧凸集之和）。

\begin{enumerate}
\def\labelenumi{\alph{enumi})}
\setcounter{enumi}{1}
\tightlist
\item
  设 E 是一个桶型空间，L 是 E 的一个子空间，它有一个具有可数基的补。证明 L 是桶型的（利用 \(a\) ) 与习题 13, b)）。
\end{enumerate}

* c) 设 E 是一个 Hausdorff 局部凸空间；它的完备化 \(\hat{\mathbf{E}}\) 可等同于一个桶型空间 F 的闭子空间，F 是一族 Fréchet 空间的乘积（II, 第 5 页, 命题 3 与 IV, 第 14 页, 推论）。设 \((e_{\alpha})_{\alpha \in \mathbf{A}}\) 是子空间 E 在 F 中的一个补的基，并设 \(\mathrm{H}_{\alpha}\) 是 F 中由 E 与诸 \(e_{\beta}\) （指标 \(\beta \neq \alpha\) ）生成的超平面；由 \(b\) )，\(\mathrm{H}_{\alpha}\) 是桶型空间。对每个 \(x \in \mathrm{E}\) ，令 \(u(x)\) 为桶型空间 \(\mathrm{G} = \prod_{\alpha \in \mathbf{A}} \mathrm{H}_{\alpha}\) （

\(F^{A}\) 的子空间）中所有坐标都等于 x 的点；u 是从 E 到子空间 \(\Delta \cap G\) 上的同构，其中 \(\Delta\) 是 \(F^{A}\) 中的对角线。证明 \(u(E) = \Delta \cap G\) 在 G 中闭，从而每个 Hausdorff 局部凸空间都同构于一个 Hausdorff 桶型空间的闭子空间。*

\begin{enumerate}
\def\labelenumi{\arabic{enumi})}
\setcounter{enumi}{14}
\tightlist
\item
  设 E 是一个 Hausdorff 桶型（相应地，次桶型）空间，\(\hat{E}\) 是它的完备化。证明：\(\hat{E}\) 的每个包含 E 的子空间 F 都是桶型的（相应地，次桶型的）（cf.~III, 第 24 页, 推论 与 IV, 第 52 页, 习题 1）。
\end{enumerate}

¶ * 16) 设 \((E_{\iota})_{\iota\in I}\) 是不可数族 Hausdorff 桶型空间，其中没有一个是点 0，并设 \(E = \prod_{\iota \in I} E_{\iota}\) ；那么 \(E\) 是桶型的（IV, 第 14 页, 推论）。设 \(G\) 是 E 中由所有点 \((x_{\iota})\) （除可数多个指标外均有 \(x_{\iota}=0\) ）组成的子空间。G 中每条在 E 中收敛的点列的极限都属于 G，但 G 在 E 中稠密。

\begin{enumerate}
\def\labelenumi{\alph{enumi})}
\tightlist
\item
  证明：每个子集 \(\mathbf{M}\) （\(\mathrm{G}' = \mathrm{E}'\) 的）若对 \(\sigma(\mathrm{E}', \mathrm{G})\) 有界，都含于一个有限乘积 \(\prod_{\iota \in \mathrm{H}} \mathrm{E}_{\iota}'\) （IV, 第 12 页, 命题 13）之中，其中 \(\mathrm{H}\) 是 I 的有限子集；因而 \(\mathbf{M}\)
\end{enumerate}

对 \(\sigma(E', E)\) 有界。由此推出 \(G\) 是桶型的。

\begin{enumerate}
\def\labelenumi{\alph{enumi})}
\setcounter{enumi}{1}
\tightlist
\item
  设 F 是 E 的一个子空间，使得 \(G \subset F \subset E\) ，且 G 是 F 中的一个（处处稠密的）超平面；F 是桶型的（习题 15）。证明 F 不是有界型的。（用反证法；若 F 中存在一个凸的、均衡的且有界的集 A，使得 G 是赋范空间 \(F_{A}\) （III, 第 7 页）中的处处稠密超平面，则将存在 G 的一个点列收敛于 F 中不属于 G 的一点）（cf.~IV, 第 52 页, 习题 2）。*
\end{enumerate}

\begin{enumerate}
\def\labelenumi{\arabic{enumi})}
\setcounter{enumi}{16}
\tightlist
\item
  \begin{enumerate}
  \def\labelenumii{\alph{enumii})}
  \tightlist
  \item
    设 E 是一个 Hausdorff 局部凸空间，L 是 E 的一个有限余维数的向量子空间，T 是 L 中一个吸有界集的桶。证明：存在 E 中一个吸有界集的桶 \(T'\) ，使得 \(T' \cap L = T\) 。（化归到 L 是 E 中超平面的情形。设 \(E_{0}\) 是与 E 相伴的有界型空间（III, 第 40 页, 习题 1），\(L_{0}\) 是赋以 \(E_{0}\) 的拓扑所诱导拓扑的超平面 L；注意 T 是 \(L_{0}\) 中 0 的一个邻域，并考虑以下两种情形：\(L_{0}\) 在 \(E_{0}\) 中稠密，或在 \(E_{0}\) 中闭；证明对 \(T'\) 可取 T 在 E 中的闭包 \(\overline{T}\) ，或 T 与一个 1 维紧凸集之和）。
  \end{enumerate}
\end{enumerate}

\begin{enumerate}
\def\labelenumi{\alph{enumi})}
\setcounter{enumi}{1}
\tightlist
\item
  设 E 是一个次桶型空间，L 是 E 的一个有限余维数的向量子空间。由 a) 推出 L 是次桶型的（cf.~IV, 第 64 页, 习题 11）。
\end{enumerate}

¶ 18) 设 E 是局部凸可度量化子空间的递增序列 ( \(E_{n}\) ) 的严格归纳极限空间（II, 第 33 页），并设 F 是 E 的一个向量子空间，使得 E 的每一点都是 F 的某个点列的极限点。

\begin{enumerate}
\def\labelenumi{\alph{enumi})}
\item
  若 \(\overline{\mathbf{E}}_n\) 是 \(\mathbf{E}_n\) 在 \(\mathbf{E}\) 中的闭包，则 \(\mathbf{E}\) 是序列 \((\overline{\mathbf{E}}_n)\) 的严格归纳极限。设 \(\mathbf{F}_n\) 是 \(\mathbf{F} \cap \overline{\mathbf{E}}_n\) 在 \(\mathbf{E}\) 中的闭包。证明 \(\mathbf{E}\) 是子空间的递增序列 \(\mathbf{F}_n\) 的并。
\item
  假设 E 是桶型的。证明 F 是有界型的。（设 u 是从 F 到一个 Banach 空间 G 中的线性映射，它把 F 的每个有界子集变为 G 的有界子集。证明存在从 E 到 G 中的一个线性映射，它在 F 上的限制等于 u，且在每个 \(F_{n}\) 上的限制连续。最后利用 III, 第 44 页, 习题 13, b)。）
\end{enumerate}

\begin{enumerate}
\def\labelenumi{\arabic{enumi})}
\setcounter{enumi}{18}
\tightlist
\item
  称 Hausdorff 局部凸空间 E 是超有界型的，如果 E 的每个吸收 E 的所有凸的、均衡的、有界的且半完备的子集的凸子集，都是 E 中 0 的邻域。
\end{enumerate}

\begin{enumerate}
\def\labelenumi{\alph{enumi})}
\item
  证明：每个超有界型空间都既是有界型的又是桶型的。
\item
  设 E 是一个 Hausdorff 局部凸空间，使得每个趋于 0 的序列的点所成之集的闭均衡凸包都是半完备的。证明：若 E 是有界型的，则它是超有界型的。特别地，每个有界型的拟完备空间都是超有界型的；每个 Fréchet 空间都是超有界型的。
\item
  设 \((\mathrm{E}_{\alpha})\) 是向量空间 E 的向量子空间的一个有向递增族，使得 E 是诸 \(\mathrm{E}_{\alpha}\) 的并。设 \(\mathcal{T}_{\alpha}\) 是 \(\mathrm{E}_{\alpha}\) 上的一个局部凸拓扑（对每个 \(\alpha\) ），并设 \(\mathcal{T}\) 是使从 \(\mathrm{E}_{\alpha}\) 到 E 中的诸典范单射都连续的最细局部凸拓扑。假设 \(\mathcal{T}\) 是 Hausdorff 的，且对每个 \(\alpha\) ，\(\mathrm{E}_{\alpha}\) 上由 \(\mathcal{T}\) 诱导的拓扑是 \(\mathcal{T}_{\alpha}\) 。证明：若诸空间 \(\mathrm{E}_{\alpha}\) 中的每一个都是超有界型的，则赋以 \(\mathcal{T}\) 的 E 是超有界型的。d) 证明：超有界型空间的每个有限乘积都是超有界型的；由此推出超有界型空间的每个拓扑直和都是超有界型的。
\item
  证明：超有界型空间的无穷序列的乘积空间 \(\mathrm{E} = \prod_{n=0}^{\infty} \mathrm{E}_n\)
\end{enumerate}

是超有界型的。（设 A 是 E 的一个凸子集，它吸收 E 的每个凸的、均衡的、有界的且半完备的子集。证明：若 A 不是 E 中 0 的邻域，则将存在一个序列 \((x_{n})\) （在 \(\mathbb{C}\) A 中），使得 \(x_{n}\) 的前 \(n - 1\) 个坐标为零，但 \(\neq 0\) 。然后注意，这样一个序列的点所成之集的闭均衡凸包等同于点 \(\sum_{n = 0}^{\infty}\lambda_n x_n\) （其中 \(\sum_{n = 0}^{\infty}|\lambda_n| \leqslant 1\) ）所成的集，且

这个包是一个半完备的集。）

\begin{enumerate}
\def\labelenumi{\arabic{enumi})}
\setcounter{enumi}{19}
\item
  证明：要使 Hausdorff 局部凸空间 \(E\) 是超有界型的，必须且只需它是 Banach 空间的一个族的归纳极限。（为看条件必要，考虑凸的、均衡的、有界的且半完备的集 \(B\) （在 \(E\) 中）以及空间 \(E_B\) 。为看条件充分，注意：若 \(E\) 是 Banach 空间的一个族 \(E_\alpha\) 的归纳极限，则可以假设诸 \(E_\alpha\) 是 \(E\) 的（代数意义上的）子空间；若 \(V\) 是 \(E\) 中吸收 \(E\) 的凸的、均衡的、有界的且半完备的子集的一个凸集，证明 \(V\) 吸收每个球 \(B_\alpha\) （\(E_\alpha\) 的）（用反证法）；若 \(V\) 不吸收 \(B_\alpha\) ，则它不吸收一个点列 \((x_n)\) （由 \(B_\alpha\) 的点组成，在 \(E\) 中趋于 0）；然后利用如下事实：在 Banach 空间中，紧集的闭凸包是紧的。）
\item
  证明：若 E 是 Hausdorff 局部凸半完备空间，则与 E 相伴的有界型空间（III, 第 40 页, 习题 1）是超有界型的。
\end{enumerate}

¶ 22) 设 E 是一个满足第一可数公理的无穷维 Banach 空间。

\begin{enumerate}
\def\labelenumi{\alph{enumi})}
\item
  证明：集 \(\mathcal{K}\) （由 E 中使 \(\mathrm{E_A}\) 为无穷维的所有紧凸均衡子集 A 组成）是无穷的，且其基数 \(\leqslant 2^{\mathrm{Card}(\mathbf{N})}\) （GT, IX, § 5, 习题 17）。对每个 \(x_0 \in \mathrm{E}\) 及每个 \(\mathrm{A} \in \mathcal{K}\) ，集 \(x_0 + \mathrm{A}\) 都含有一个基数为 \(2^{\mathrm{Card}(\mathbf{N})}\) 的自由子集（II, 第 80 页, 习题 24, c)）。
\item
  设 E 中 \(x_0 \neq 0\) 。证明：存在一个族 \((y_{\mathrm{A}})_{\mathrm{A} \in \mathcal{K}}\) ，使得 \(x_0\) 与诸 \(y_{\mathrm{A}}\) 构成一个自由族，且有 \(y_{\mathrm{A}} \in x_0 + \mathrm{A}\) （对一切 \(\mathrm{A} \in \mathcal{K}\) ；将 \(\mathcal{K}\) 良序化，并利用 \(a\) ) 作超限归纳论证）。
\item
  设 \(f \in \mathrm{E}^{*}\) 是满足 \(f(x_0) = 1\) 且 \(f(y_{\mathrm{A}}) = 0\) （对一切 \(\mathrm{A} \in \mathcal{K}\)）的线性型，并设 \(\mathrm{H} = f^{-1}(0)\) 。证明：一个子集 \(\mathbf{M}\) （\(\mathrm{H}\) 的）若是凸的、均衡的且半完备的，则必是有限维的（注意否则 \(\mathbf{M}\) 将含有一个无穷维的紧凸均衡集 \(\mathrm{A}\) ；于是 \(y_{\mathrm{A}}\) 将属于 \(\mathrm{H} \cap (x_0 + \mathbf{M})\) ）。
\item
  证明：赋以 E 的拓扑所诱导拓扑的 H 不是超有界型的，尽管它是有界型的且是桶型的（III, 第 44 页, 习题 14）。（利用 c)，证明若 H 是超有界型的，则其拓扑将是最细局部凸拓扑，由此得出矛盾。）
\end{enumerate}

\exercisegroup

\begin{enumerate}
\def\labelenumi{\arabic{enumi})}
\item
  设 E、F 与 G 是三个局部凸空间，S 是 E 的由有界集组成的一个覆盖。证明：若 u 是从 \(E \times F\) 到 G 中的分别连续双线性映射，使得对每个集 \(M \in S\) ，u 在 \(M \times F\) 上的限制连续，则 u 是 S-亚连续的。
\item
  设 \(\mathbf{E}\) 是直和空间 \(\mathbf{R}^{(\mathbf{N})}\) ，赋以 \(\mathbf{R}^{\mathbf{N}}\) 上乘积拓扑所诱导的拓扑。证明：双线性型 \(((x_n), (y_n)) \mapsto \sum_{n=0}^{\infty} x_n y_n\) （在 \(\mathbf{E} \times \mathbf{E}\) 上）是分别连续的，但对每个集 \(\mathfrak{S}\) （由 \(\mathbf{E}\) 的有界子集组成，且至少含有一个无穷维有界集），这个双线性型都不是 \(\mathfrak{S}\) -亚连续的。
\end{enumerate}


\begin{enumerate}
\def\labelenumi{\arabic{enumi})}
\setcounter{enumi}{2}
\tightlist
\item
  设 E 是赋以最细局部凸拓扑（II, 第 26 页）的空间 \(\mathbf{R}^{(\mathbf{N})}\) ；设 F 是空间 \(R^{N}\) ；空间 E 是超有界型的（III, 第 45 页, 习题 19）且完备，而 F 可度量化且完备。设 S（相应地，T）是 E（相应地，F）的所有有界子集组成的集。证明：E × F 上的双线性型 \(((x_{n}), (y_{n})) \mapsto \sum_{n=0}^{\infty} x_{n} y_{n}\) 是 \((\mathfrak{S}, \mathfrak{T})\) -亚连续的，但不连续
\end{enumerate}

（cf.~IV, 第 48 页, 习题 11）。

\begin{enumerate}
\def\labelenumi{\arabic{enumi})}
\setcounter{enumi}{3}
\item
  设 E 是一个局部凸空间，F 是一个次桶型空间（III, 第 44 页, 习题 7），T 是 F 的所有有界子集组成的集。证明：若从 \(E \times F\) 到一个局部凸空间 G 中的双线性映射是 T-亚连续的，则对 E 的有界子集的每个集 S，它都是 \((\mathfrak{S}, \mathfrak{T})\) -亚连续的（cf.~III, 第 44 页, 习题 11）。
\item
  \begin{enumerate}
  \def\labelenumii{\alph{enumii})}
  \tightlist
  \item
    设 \(\mathbf{E}\) 、\(\mathbf{F}\) 与 \(\mathbf{G}\) 是三个 Hausdorff 局部凸空间，\(u\) 是从 \(\mathrm{E} \times \mathrm{F}\) 到 \(\mathbf{G}\) 中的一个双线性映射。要使存在 0 的一个均衡邻域 \(\mathbf{U}\) （\(\mathbf{E}\) 中），使得映射 \(u(x, .)\) （当 \(x\) 取遍 \(\mathbf{U}\) 时）全体所成的集在 \(\mathcal{L}(\mathbf{F}; \mathbf{G})\) 中等度连续，必须且只需：\(u\) 当把 \(\mathbf{E}\) 的拓扑换成使诸集 \(\lambda \mathbf{U} (\lambda \neq 0)\) 构成 0 的基本邻域系的最粗拓扑时是连续的。证明：若 \(\mathbf{G}\) 赋范，则从 \(\mathrm{E} \times \mathrm{F}\) 到 \(\mathbf{G}\) 中的每个连续双线性映射都满足这一条件。
  \end{enumerate}
\end{enumerate}

\begin{enumerate}
\def\labelenumi{\alph{enumi})}
\setcounter{enumi}{1}
\tightlist
\item
  取乘积空间 \(R^{N}\) 作为 E、F 与 G，并取连续双线性映射 \(((x_{n}), (y_{n})) \mapsto (x_{n}y_{n})\) 作为 u。证明：不存在 E 中 0 的任何邻域 U，使得当 x 取遍 U 时的映射 \(u(x, .)\) 全体所成的集在 \(\mathcal{L}(F; G)\) 中等度连续。
\end{enumerate}

\begin{enumerate}
\def\labelenumi{\arabic{enumi})}
\setcounter{enumi}{5}
\tightlist
\item
  设 E、F 与 G 是三个拓扑向量空间。称从 \(E \times F\) 到 G 中的双线性映射的一个集 H 是分别等度连续的，如果对一切 \(x \in E\) ，当 u 取遍 H 时的线性映射 \(u(x, .)\) 全体所成的集在 \(\mathcal{L}(F; G)\) 中等度连续，且对一切 \(y \in F\) ，当 u 取遍 H 时的线性映射 \(u(. , y)\) 全体所成的集在 \(\mathcal{L}(E; G)\) 中等度连续。
\end{enumerate}

假设 F 可度量化，且 E 是 Baire 空间（cf.~III, 第 43 页, 习题 5 与 V, 第 79 页, 习题 15）。证明：从 \(E \times F\) 到 G 中的双线性映射的每个分别等度连续的集都是等度连续的（cf.~III, 第 42 页, 习题 11）。

\begin{enumerate}
\def\labelenumi{\arabic{enumi})}
\setcounter{enumi}{6}
\tightlist
\item
  设 \(\mathrm{E}\) 、\(\mathrm{F}\) 与 \(\mathrm{G}\) 是三个拓扑向量空间，\(\mathfrak{S}\) 是 \(\mathrm{E}\) 的有界子集的一个集，\(\mathrm{H}\) 是从 \(\mathrm{E} \times \mathrm{F}\) 到 \(\mathrm{G}\) 中的分别连续双线性映射的一个集。下列性质等价：
\end{enumerate}

\(\alpha)\) 对 0 的每个邻域 \(\mathbf{W}\) （\(\mathbf{G}\) 中）以及每个集 \(\mathbf{M} \in \mathfrak{S}\) ，存在 0 的一个邻域 \(\mathbf{V}\) （\(\mathbf{F}\) 中），使得 \(u(\mathbf{M} \times \mathbf{V}) \subset \mathbf{W}\) （对一切 \(u \in \mathbf{H}\)）。

\(\beta)\) 对每个集 \(\mathbf{M} \in \mathfrak{S}\) ，\(\mathbf{H} \times \mathbf{M}\) 在映射 \((u, x) \mapsto u(x, .)\) 下的像是 \(\mathcal{L}(\mathrm{F}; \mathrm{G})\) 的一个等度连续子集。

\(\gamma)\) 当 \(u\) 取遍 H 时，映射 \(y \mapsto u(., y)\) （从 F 到 \(\mathcal{L}_{\mathfrak{S}}(\mathrm{E}; \mathrm{G})\) 中）全体所成的集是等度连续的。

这时称 H 是从 \(E \times F\) 到 G 中的（分别连续的）双线性映射的一个 S-等度亚连续集。类似地，对 F 的有界子集的一个集 T，我们定义 T-等度亚连续集与 \((\mathfrak{S}, \mathfrak{T})\) -等度亚连续集的概念。

\begin{enumerate}
\def\labelenumi{\arabic{enumi})}
\setcounter{enumi}{7}
\item
  设 H 是从 \(E \times F\) 到 G 中的双线性映射的一个 S-等度亚连续集（习题 7）。对每个子集 \(M \in S\) ，证明 H 在 \(M \times F\) 上等度连续；此外，对 F 的每个有界子集 Q，诸集 \(u(\mathbf{M} \times \mathbf{Q})\) （其中 u 取遍 H）的并在 G 中有界。
\item
  设 H 是双线性映射的一个 \((\mathfrak{S}, \mathfrak{T})\) -等度亚连续集（从 \(E \times F\) 到 G 中）（习题 7）；证明：对每一对集 \(M \in S\) 、\(N \in T\) ，H 在 \(M \times N\) 上一致等度连续。
\item
  设 \(\mathrm{E}_1\) 、\(\mathrm{E}_2\) 与 \(\mathrm{F}\) 是三个拓扑向量空间，\(\mathbf{G}_1\) （相应地，\(\mathbf{G}_2\) ）是 \(\mathrm{E}_1\) （相应地，\(\mathrm{E}_2\) ）的一个处处稠密子空间，\(\mathfrak{S}_1\) （相应地，\(\mathfrak{S}_2\) ）是 \(\mathbf{G}_1\) （相应地，\(\mathbf{G}_2\) ）的有界子集的一个族。设 \(\mathrm{H}\) 是从 \(\mathrm{E}_1 \times \mathrm{E}_2\) 到 \(\mathrm{F}\) 中的分别连续双线性映射的一个集；若在 \(\mathbf{G}_1 \times \mathbf{G}_2\) 上取诸映射 \(u \in \mathrm{H}\) 的限制，所得之集是 \((\mathfrak{S}_1, \mathfrak{S}_2)\) -等度亚连续的，则 \(\mathrm{H}\) 亦是。
\item
  若 \(\mathbf{F}\) 是桶型空间，则从 \(\mathbf{E} \times \mathbf{F}\) 到一个局部凸空间 \(\mathbf{G}\) 中的双线性映射的每个分别等度连续的集都是 \(\mathfrak{S}\) -等度亚连续的（对每个集 \(\mathfrak{S}\) ，由 \(\mathbf{E}\) 的有界子集组成）。
\item
  设 E、F 是两个拓扑向量空间，f 是双线性映射 \((x, u) \mapsto u(x)\) （从 \(E \times \mathcal{L}(E; F)\) 到 F 中）；设 T 是 \(\mathcal{L}(E; F)\) 上与向量空间结构相容且比简单收敛拓扑更细的一个拓扑。设 S 是 E 的有界子集的一个族，U 是 \(\mathcal{L}(E; F)\) 的（对拓扑 T 而言的）有界子集的一个族。证明：f 是 S-亚连续的，当且仅当 T 比 S-拓扑细；f 是 U-亚连续的，当且仅当 U 中诸集是 \(\mathcal{L}(E; F)\) 的等度连续子集。
\item
  设 E、F、G 是三个拓扑向量空间，S（相应地，T）是 E（相应地，F）的有界子集的一个族。设 H 是从 \(E \times F\) 到 G 中的 T-亚连续双线性映射所组成的向量空间。
\end{enumerate}

\begin{enumerate}
\def\labelenumi{\alph{enumi})}
\item
  证明：在 H 上，对形如 \(M \times N\) （其中 \(M \in S\) 、\(N \in T\) ）的集合的一致收敛拓扑与向量空间结构相容；这个拓扑称为 H 上的 \((\mathfrak{S}, \mathfrak{T})\) -拓扑。对每个映射 \(u \in H\) ，令 \(\tilde{u}\) 为连续映射 \(x \mapsto u(x, .)\) （从 E 到 \(\mathcal{L}_{\mathfrak{T}}(F; G)\) 中）。证明 \(u \mapsto \tilde{u}\) 是从赋以 \((\mathfrak{S}, \mathfrak{T})\) -拓扑的空间 H 到空间 \(\mathcal{L}_{\mathfrak{S}}(E; \mathcal{L}_{\mathfrak{T}}(F; G))\) 上的一个同构。
\item
  设 L 是 H 的一个子集，使得对每一对 \((x, y) \in E \times F\) ，当 u 取遍 L 时的 \(u(x, y)\) 全体所成的集在 G 中有界（H 的简单有界子集）。证明：若 E、F、G 是局部凸的，且 E 与 F 是 Hausdorff 的和拟完备的，则 L 在 H 中对 \((\mathfrak{S}, \mathfrak{T})\) -拓扑有界。
\item
  设 E、F、G 是三个 Hausdorff 局部凸空间。若 E 是桶型的且 F 拟完备或桶型，则 H 的每个简单有界子集 L 都是 T-等度亚连续的（III, 第 47 页, 习题 7）。
\end{enumerate}

\(d\) ) 若 E 与 F 是桶型的，G 拟完备，且 \(\mathfrak{S}\) 与 \(\mathfrak{T}\) 分别是 E 与 F 的覆盖，则 H 对 \((\mathfrak{S},\mathfrak{T})\) -拓扑是 Hausdorff 的且拟完备的。

\begin{enumerate}
\def\labelenumi{\arabic{enumi})}
\setcounter{enumi}{13}
\item
  把 § 5 的定义与结果推广到任意多重线性映射。设 E、F、G 是三个拓扑向量空间，S（相应地，T）是 E（相应地，F）的有界子集的一个族，U 是从 E × F 到 G 中的双线性 (S, T)-亚连续映射所组成的空间 \(\mathcal{L}_{\mathfrak{S},\mathfrak{T}}\) (E, F; G)（赋以 (S, T)-拓扑（习题 13））的有界子集的一个族。证明：从 E × F × \(\mathcal{L}_{\mathfrak{S},\mathfrak{T}}\) (E, F; G) 到 G 中的三线性映射 (x, y, u) → u(x, y) 是 (S, T)-亚连续的；要使它是 (S, U)-亚连续的，必须且只需 U 中的每个集 L 都是 S-等度亚连续的（III, 第 47 页, 习题 7）。
\item
  设 \(E\) 是由所有实数序列 \(x = (\xi_n)_{n \geqslant 0}\) （使通项为 \(\xi_n\) 的级数收敛）组成的空间。令 \(\| x \| = \sup_n | \sum_{k=0}^n \xi_k |\) 。
\end{enumerate}

\begin{enumerate}
\def\labelenumi{\alph{enumi})}
\tightlist
\item
  证明 \(\| x\|\) 是 E 上的一个范数，且 E 对这个范数是完备的。
\end{enumerate}

\begin{enumerate}
\def\labelenumi{\alph{enumi})}
\setcounter{enumi}{1}
\tightlist
\item
  证明：向量空间 \(\ell^1 (\mathbf{N})\) （I, 第 4 页）作为 E 的子空间（对 E 的拓扑而言）是处处稠密的；\(\ell^1 (\mathbf{N})\) 上由范数 \(\| x\| _1 = \sum_{n = 0}^{\infty}|\xi_n|\) 定义的拓扑严格地比 E 的拓扑所诱导的拓扑细。
\end{enumerate}

\begin{enumerate}
\def\labelenumi{\alph{enumi})}
\setcounter{enumi}{2}
\tightlist
\item
  设 \((\mathbf{P}_n)\) 是 \(\mathbf{N} \times \mathbf{N}\) 的有限子集的递增序列，构成 \(\mathbf{N} \times \mathbf{N}\) 的一个覆盖。对每个 \(x = (\xi_n) \in \mathrm{E}\) 及每个 \(y = (\eta_n) \in \ell^1(\mathbf{N})\) ，令 \(f_n(x, y) = \sum_{(i,j) \in \mathbf{P}_n} \xi_i \eta_j\) 。那么，序列 \((f_n(x, y))\) 对每对 \((x, y) \in \mathrm{E} \times \ell^1(\mathbf{N})\) 都趋于一个极限，当且仅当对这些对 \((x, y)\) 中的每一对，序列 \((f_n(x, y))\) 有界；此时 \(f_n(x, y)\) 的极限等于 \(\left(\sum_{n=0}^{\infty} \xi_n\right) \left(\sum_{n=0}^{\infty} \eta_n\right)\) 。
\end{enumerate}

（利用 III, 第 48 页, 习题 13, c)，证明双线性型的序列 \((f_n)\) 等度连续，并注意它在子空间 \(\ell^1 (\mathbf{N})\times \ell^1 (\mathbf{N})\) 中收敛；再利用 \(b).\) 完成论证。）\(d)\) 对每个 \(j\in \mathbf{N}\) ，令 \(\rho_{jn}\) 为 \(\mathbf{N}\) 的闭区间的最小个数，使这些闭区间的并是 \(\mathbf{P}_n\cap (\mathbf{N}^{\prime}\times \{j\})\) 在 \(\mathbf{N}\) 上的投影；令 \(\rho_{n} = \sup_{j\in \mathbf{N}}\rho_{jn}\) 。证明在 \(c)\) 中得到的条件等价于 \(\sup_n\rho_n < + \infty\) 。（若 \(\phi_n\) 是 \(\mathbf{P}_n\) 的特征函数，证明

双线性型 \(f_{n}\) 的范数是 \(\sup_{j\in \mathbf{N}}\big(\sum_{i = 0}^{\infty}|\rho_n(i,j) - \phi_n(i + 1,j)|\big).\)

\exercisegroup

¶ 1) Hausdorff 局部凸空间的一个穷竭由一个筛 \(C = (C_n, p_n)_{n \geqslant 0}\) （GT, IX, § 6, No.~5, 定义 8）以及对每个 \(n \geqslant 0\) 、一个映射 \(\rho_n\) （从 \(C_n\) 到 E 的所有凸均衡子集所成的集中）给出，它具有下列性质：E1) E 是诸 \(\phi_0(c)\) （其中 \(c\) 取遍 \(C_0\) ）的并；

E2) 对每个 \(n\) 及每个 \(c \in \mathbf{C}_n\) ，\(\phi_n(c)\) 是诸 \(\phi_{n+1}(c')\) （其中 \(c'\) 取遍 \(p_n^{-1}(c)\) ）的并；E3) 对每个序列 \((c_k)_{k \geqslant 0}\) （满足 \(c_k \in \mathbf{C}_k\) 且 \(c_k = p_k(c_{k+1})\) ，对一切 \(k \geqslant 0\) ），存在一个序列 \((\phi_k)\) （由 \(>0\) 的数组成），使得对每个序列 \((x_k)\) （E 的点列，满足 \(x_k \in \phi_k(c_k)\) ）以及每个序列 \((\lambda_k)\) （实数的，满足 \(0 \leqslant \lambda_k \leqslant \rho_k\) ，对一切 \(k\) ），级数 \(\sum_{k=0}^{\infty} \lambda_k x_k\) 在 E 中收敛。

\begin{enumerate}
\def\labelenumi{\alph{enumi})}
\tightlist
\item
  在上述假设下，证明：若此外诸 \(\phi_n(c)\) （对 \(c \in \mathbf{C}_n\) ）还是闭的，则可以假设诸 \(\rho_k\) 已如此选取，使得 \(\sum_{k=m}^{\infty} \lambda_k x_k \in \phi_m(c_m)\) （对一切 \(m \geqslant 1\)）。

（取诸 \(\rho_{k}\) 使得 \(\sum_{k=0}^{\infty} \rho_{k} \leqslant 1\) ）。
\end{enumerate}

\begin{enumerate}
\def\labelenumi{\alph{enumi})}
\setcounter{enumi}{1}
\tightlist
\item
  假设给定一个筛 C 以及映到 E 的凸均衡子集所成之集中的映射的序列 \((\phi_{n})\) ，满足 E1)、E2) 以及下列条件：E3') 对每个序列 \((c_{k})_{k\geq0}\) （满足 \(c_{k}=p_{k}(c_{k+1})\) ，对一切 \(k\geq0\) ），存在一个序列 \((\mu_{k})\) （由 \textgreater0 的数组成），使得对每个序列 \((x_{k})\) （E 的点列，满足 \(x_{k}\in\phi_{k}(c_{k})\) ，对一切 k），点列 \((\mu_{k}x_{k})\) 含于 E 中一个凸的、有界均衡的且半完备的集中。
\end{enumerate}

证明此时条件 E3) 也成立（取 \(\rho_{k} = 2^{-k}\mu_{k}\) ）。

称局部凸 Hausdorff 空间为可穷竭的，如果存在 E 的一个穷竭。

\begin{enumerate}
\def\labelenumi{\arabic{enumi})}
\setcounter{enumi}{1}
\tightlist
\item
  设 \(\mathbf{E}\) 是一个局部凸空间且为 Baire 空间，\(\mathbf{F}\) 是一个局部凸可穷竭空间（习题 1），\((\mathbf{C}_n, p_n, \phi_n)\) 是 \(\mathbf{F}\) 的一个穷竭。
\end{enumerate}

\begin{enumerate}
\def\labelenumi{\alph{enumi})}
\item
  设 \(u\) 是从 \(\mathbf{E}\) 到 \(\mathbf{F}\) 中的一个线性映射，\(\mathbf{W}\) 是 \(\mathbf{F}\) 中一个凸的、均衡的且吸收的集。证明：存在一个序列 \((c_k)\) ，使得 \(c_k \in \mathbf{C}_k\) ，\(c_k = p_k(c_{k+1})\) （对一切 \(k \geqslant 0\) ），以及一个整数序列 \((m_k)\) （诸整数 \(> 0\) ），使得诸集 \(u^{-1}(\phi_k(c_k) \cap m_k\mathbf{W})\) （记作 \(\mathbf{M}_k\) ）中的每一个都不是 \(\mathbf{E}\) 中的贫集。证明：对每个 \(\varepsilon > 0\) ，存在一个数列 \((\nu_k)\) （诸数 \(> 0\) ），使得若点列 \((x_k)_{k \geqslant 1}\) （\(\mathbf{E}\) 的）满足 \(x_k \in \nu_k\mathbf{M}_k\) （对一切 \(k \geqslant 1\) ），则级数 \(\sum_{k=1}^{\infty} u(x_k)\) 在 \(\mathbf{F}\) 中收敛且其和属于 \(\varepsilon\mathbf{W}\) 。
\item
  此外假设 E 可度量化，且 u 在 \(E \times F\) 中的图像是闭的。证明：对每个 \(\varepsilon > 0\) ，有 \(\overline{u^{-1}(\mathbf{W})} \subset (1 + \varepsilon) u^{-1}(\mathbf{W})\) 。（注意：若 \((\mathrm{U}_{k})\) 是 E 中 0 的可数基本邻域系，则对每个 k 存在 E 中 0 的一个凸均衡邻域 \(V_{k}\) ，使得 \(V_{k} \subset U_{k} \cap v_{k} \overline{M}_{k}\) 。对每一点 \(a \in \overline{u^{-1}(\mathbf{W})}\) ，求一个序列 \((x_{k})_{k \geqslant 0}\) ，使得 \(x_{0} \in u^{-1}(\mathbf{W})\) ，\(x_{k} \in v_{k} M_{k}\) （对 \(k \geqslant 1\) ），且 \(a - \sum_{j=0}^{k} x_{j} \in V_{k}\) （对一切 \(k \geqslant 1\) ），然后应用 a)。
\item
  由 \(b\) ) 推出：若 \(E\) 是可度量化 Baire 空间，则从 \(E\) 到 \(F\) 中的每个图像为闭的线性映射都连续。
\end{enumerate}

\begin{enumerate}
\def\labelenumi{\arabic{enumi})}
\setcounter{enumi}{2}
\item
  证明：Fréchet 空间 \(\mathbf{E}\) 是可穷竭的（若 \((\mathbf{U}_k)\) 是构成 \(\mathbf{E}\) 中 0 的闭的凸均衡邻域的基本邻域系的一个递减序列，考虑诸集 \((m + 1)\mathrm{U}_k\) （其中 \(m\) 与 \(k\) 取遍 \(\mathbf{N}\) ）的有限交）。
\item
  \begin{enumerate}
  \def\labelenumii{\alph{enumii})}
  \tightlist
  \item
    可穷竭局部凸空间的每个闭子空间都是可穷竭的。
  \end{enumerate}
\end{enumerate}

\begin{enumerate}
\def\labelenumi{\alph{enumi})}
\setcounter{enumi}{1}
\tightlist
\item
  设 E 是一个可穷竭局部凸空间，\(u:E \rightarrow F\) 是从 E 到 F 上的一个连续线性满射映射。证明 F 是可穷竭的。特别地，E 对其闭子空间的每个商空间都是可穷竭的。在 E 上赋以比 E 的拓扑粗的 Hausdorff 局部凸拓扑所得到的每个空间都是可穷竭的。
\end{enumerate}

\begin{enumerate}
\def\labelenumi{\arabic{enumi})}
\setcounter{enumi}{4}
\tightlist
\item
  设 \((\mathrm{C}^{(m)})_{m\geqslant 0}\) 是筛 \(\mathrm{C}^{(m)} = (\mathrm{C}_n^{(m)},p_n^{(m)})_{n\geqslant 0}\) 的一个序列。对每个 \(n\geqslant 0\) ，令
\end{enumerate}

\[
\mathrm{D} _ {n} = \mathrm{C} _ {n} ^ {(0)} \times \mathrm{C} _ {n - 1} ^ {(1)} \times \dots \times \mathrm{C} _ {0} ^ {(n)} \times \prod_ {m = n + 1} ^ {\infty} \left\{a _ {m} \right\},
\]

其中 \(a_{m} = 0\) （对一切 \(m\geqslant 0\) ）；映射 \(p_n:\mathrm{D}_{n + 1}\to \mathrm{D}_n\) 取为等于

\[
p _ {n} ^ {(0)} \times p _ {n - 1} ^ {(1)} \times \dots \times p _ {0} ^ {(n)} \times q ^ {(n + 1)} \times \prod_ {m = n + 2} ^ {\infty} \mathrm{id} _ {m}
\]

其中 \(q^{(n + 1)}\) 是从 \(C_0^{(n + 1)}\) 到 \(\{0\}\) 上的唯一映射，\(\mathrm{id}_m\) 是 \(\{a_m\}\) 的恒等映射。那么 \((\mathrm{D}_n, p_n)\) 是一个筛。

设 \((\mathrm{E}^{(m)})_{m\geqslant 0}\) 是 Hausdorff 局部凸空间的一个序列；我们假设对每个 \(m\) 存在一个穷竭 \((\mathrm{C}_n^{(m)},p_n^{(m)},\phi_n^{(m)})_{n\geqslant 0}\) （\(\mathrm{E}^{(m)}\) 的）。考虑 Hausdorff 局部凸空间 \(\mathrm{E} = \prod \mathrm{E}^{(m)}\) ，且对每个 \(n\) ，令

\[
\phi_ {n} = \phi_ {n} ^ {(0)} \times \phi_ {n - 1} ^ {(1)} \times \dots \times \phi_ {0} ^ {(n)} \times \prod_ {m = n + 1} ^ {\infty} \psi_ {m}
\]

其中 \(\psi_{m}\) 是从 \(\{a_m\}\) 到 \(\mathrm{E}^{(m)}\) 的凸均衡子集所成之集中的映射，使得 \(\psi_{n}(a_{m}) = \mathrm{E}^{(m)}\) 。证明 \((\mathrm{D}_n,p_n,\phi_n)\) 是乘积空间 E 上的一个穷竭。

\begin{enumerate}
\def\labelenumi{\arabic{enumi})}
\setcounter{enumi}{5}
\tightlist
\item
  证明：递增的子空间序列 \(\mathbf{E}_n\) （向量空间 \(\mathbf{E}\) 的）的归纳极限（II, 第 31 页）（赋以拓扑 \(\mathcal{T}_n\) ，且 \(\mathbf{E}_n\) 赋以 \(\mathcal{T}_n\) 后可穷竭）若为 Hausdorff 的，则是一个可穷竭局部凸空间。
\end{enumerate}

\chapter{拓扑向量空间中的对偶}

本章中，所考虑的所有向量空间都是域 K 上的向量空间，其中 K 为 R 或 C。

\section{对偶}

\subsection{与对偶相容的拓扑}

本节中，E 与 F 表示由一个双线性型 B （II, 第 40 页）置于对偶的两个向量空间。我们回顾（II, 第 41 页）曾定义过两个线性映射

\[
d _ {\mathbf {B}}: \mathrm{F} \to \mathrm{E} ^ {*}, s _ {\mathbf {B}}: \mathrm{E} \to \mathrm{F} ^ {*}
\]

它们由关系式

\[
\mathbf {B} (x, y) = \langle x, d _ {\mathbf {B}} (y) \rangle = \langle y, s _ {\mathbf {B}} (x) \rangle\tag{1}
\]

刻画，其中 \(x\in \mathbf{E},y\in \mathbf{F}.\)

定义 1. --- E 上的局部凸拓扑 T 称为与 E 和 F 之间的对偶相容，如果 \(d_{B}\) 是从 F 到由赋以拓扑 T 的 E 所得局部凸空间的对偶上的一个双射。

如果存在这样一个拓扑 \(\mathcal{T}\) ，则映射 \(d_{\mathrm{B}}\) 是单射，也就是说，E 与 F 之间的对偶在 F 中是分离的（II, 第 41 页）。

命题 1. --- (i) E 中的闭凸子集对 E 上所有与 E 和 F 之间的对偶相容的局部凸拓扑来说是相同的。

\begin{enumerate}
\def\labelenumi{(\roman{enumi})}
\setcounter{enumi}{1}
\tightlist
\item
  \(\mathbf{E}\) 的有界子集对 \(\mathbf{E}\) 上所有与 \(\mathbf{E}\) 和 \(\mathbf{F}\) 之间的对偶相容的局部凸拓扑来说是相同的。
\end{enumerate}

设 T 是 E 上与 E 和 F 之间的对偶相容的一个拓扑，因而比 \(\sigma(\mathrm{E},\mathrm{F})\) 细。若 E 的一个凸子集对 T 是闭的，则它是闭的实半空间之交（II, 第 38 页, 推论 1），因而它对 \(\sigma(\mathrm{E},\mathrm{F})\) 是闭的。这就证明了 (i)。断言 (ii) 已在 III, 第 27 页, 推论 3 中证明。

用 \(F_{\sigma}\) 表示赋以弱拓扑 \(\sigma(F, E)\) 的向量空间 F。那么线性映射 \(s_{B}\) 把 E 映到对偶 \((F_{\sigma})'\) （\(F_{\sigma}\) 的）上（II, 第 43 页, 命题 3）。设 \(S\) 是 \(F_{\sigma}\) 的有界子集的一个族。滥用术语，我们把在 \(s_{B}\) 之下的 \(S\) -拓扑（\((F_{\sigma})'\) 上的）的原像称为 E 上的 \(S\) -拓扑。它由半范数族

\[
p _ {\mathbf {A}} (x) = \sup _ {y \in \mathbf {A}} | \mathrm{B} (x, y) |,\tag{2}
\]

所定义，其中 A 取遍 S。特别地，当 S 是 F 的有限子集的族时，S-拓扑正是弱拓扑 \(\sigma(\mathrm{E}, \mathrm{F})\) 。

定义 2. --- 设 E 与 F 是两个处于对偶的空间。E 上的 Mackey 拓扑，记作 \(\tau(\mathrm{E},\mathrm{F})\) ，定义为 E 上的 \(\mathfrak{S}\) -拓扑，其中 \(\mathfrak{S}\) 是 F 的这样一些子集全体所成的族：它们在 \(E^{*}\) 中的像（在 \(d_{\mathrm{B}}\) 之下）对 \(\sigma(\mathrm{E}^{*},\mathrm{E})\) 是凸的、均衡的且紧的。

当 E 与 F 之间的对偶在 F 中是分离的时，\(d_{B}\) 是单射，且 F 上的拓扑 \(\sigma(F, E)\) 是在 \(d_{B}\) 之下的拓扑 \(\sigma(E^{*}, E)\) （\(E^{*}\) 上的）的原像。在这种情形，S 由 F 的对 \(\sigma(F, E)\) 为凸的、均衡的且紧的那些子集组成。

一般地，若 \(F_{1}=d_{\mathrm{B}}(\mathrm{F})\subset\mathrm{E}^{*}\) ，并且用 \((x,y_{1})\mapsto\mathrm{B}_{1}(x,y_{1})\) 表示典范双线性型 \((x,x^{*})\mapsto\langle x,x^{*}\rangle\) 在 \(E\times F_{1}\) 上的限制，则 E 与 \(F_{1}\) 由 \(B_{1}\) 置于对偶，且这个对偶在 \(F_{1}\) 中是分离的，因为按定义有 \(\mathrm{B}(x,y)=B_{1}(x,d_{\mathrm{B}}(y))\) ，定义 2 表明 \(\tau(\mathrm{E},\mathrm{F})=\tau(\mathrm{E},\mathrm{F}_{1})\) 。

注记 1. --- 设 A 是 Hausdorff 局部凸空间 G 的一个紧凸子集，\(\tilde{A}\) 是 A 的闭均衡凸包。当域 K 为 R 时，集 \(\tilde{A}\) 是 \(A \cup (-A)\) 的闭凸包；当 K 为 C 时，集 \(\tilde{A}\) 含于 \(2A \cup (-2A) \cup (2iA) \cup (-2iA)\) 的闭凸包。因此（II, 第 14 页, 命题 15），\(\tilde{A}\) 是紧的。

特别地，由此推出，当 E 与 F 之间的对偶在 F 中是分离的时，Mackey 拓扑 \(\tau(\mathrm{E},\mathrm{F})\) 也是 \(S'\) -拓扑，其中 \(S'\) 是 F 的对 \(\sigma(\mathrm{F},\mathrm{E})\) 为紧的所有凸子集所成的集。以类似的方式我们在 F 上定义 Mackey 拓扑 \(\tau(\mathrm{F},\mathrm{E})\) 。

定理 1（Mackey）. --- 设 E 与 F 是两个处于对偶的空间；假设这个对偶在 F 中是分离的。要使 E 上的局部凸拓扑 T 与 E 和 F 之间的对偶相容，必须且只需 T 比拓扑 \(\sigma(\mathrm{E},\mathrm{F})\) 细且比 Mackey 拓扑 \(\tau(\mathrm{E},\mathrm{F})\) 粗。

把 F 与它在 \(E^{*}\) 中的像（在 \(d_{B}\) 之下）等同。用 \(S_{0}\) 表示 F 的对 \(\sigma(F, E)\) 为凸的、均衡的且紧的所有子集所成的集。按定义，\(\tau(E, F)\) 是 E 上的 \(S_{0}\) -拓扑，因而比 \(\sigma(E, F)\) 细。

引理 1. --- E* 的子空间 F 由 E 上对 \(\tau(\mathrm{E},\mathrm{F})\) 连续的所有线性型组成。

F 的每个元素都是对 \(\sigma(\mathrm{E},\mathrm{F})\) 连续的映射，因而对 \(\tau(\mathrm{E},\mathrm{F})\) 连续。

反之，设 \(f \in \mathrm{E}^*\) 对 \(\tau(\mathrm{E}, \mathrm{F})\) 连续。存在 0 的一个邻域 \(\mathbf{U}\) （\(\mathrm{E}\) 中，对 \(\tau(\mathrm{E}, \mathrm{F})\) 而言），使得 \(|f| \leqslant 1\) 在 \(\mathbf{U}\) 上成立；可以假设存在一个集 \(\mathbf{A} \in \mathfrak{S}_0\) 使得 \(\mathbf{U} = \mathbf{A}^\circ\) 。换言之，\(f\) 属于双极集 \(\mathbf{A}^{\circ\circ}\) （\(\mathbf{A}\) 关于 \(\mathbf{E}^*\) 与 \(\mathbf{E}\) 之间的对偶而言）。但拓扑 \(\sigma(\mathrm{F}, \mathrm{E})\) （\(\mathbf{F}\) 上的）由 \(\sigma(\mathrm{E}^*, \mathrm{E})\) 诱导；于是 \(\mathbf{A}\) 对 \(\sigma(\mathrm{E}^*, \mathrm{E})\) 是凸的、均衡的且紧的，且双极定理（II, 第 44 页, 定理 1）蕴含等式 \(\mathbf{A} = \mathbf{A}^{\circ\circ}\) 。因此有 \(f \in \mathrm{F}\) ，引理由此得证。

引理 2. --- 设 T 是 E 上的一个局部凸拓扑，使得 E 上对 T 连续的每个线性型都属于 F。那么 T 比 \(\tau(\mathrm{E}, \mathrm{F})\) 粗。

设 \(\mathfrak{U}\) 是对 \(\mathcal{T}\) 而言 0 的凸均衡邻域所成的集。设 \(\mathfrak{S}\) 是 \(\mathfrak{U}\) 中诸元素在 F 中的极集所成的集。由 III, 第 17 页, 推论 2，有 \(\mathfrak{S} \subset \mathfrak{S}_0\) ；又由 III, 第 19 页, 命题 7 的推论 1，\(\mathcal{T}\) 与 \(\mathfrak{S}'\) -拓扑恒同，其中 \(\mathfrak{S}'\) 是 \(\mathfrak{U}\) 中诸集在 E 的对偶 E' 中的极集所成的集。但由假设 E' \(\subset\) F，因而 \(\mathfrak{S}'\) 的每个集都含于 \(\mathfrak{S}\) 的某个集中；引理由此得证。

设 T 是 E 上与 E 和 F 之间的对偶相容的一个拓扑。由引理 2，T 比 \(\tau(\mathrm{E},\mathrm{F})\) 粗，且显然 T 比 \(\sigma(\mathrm{E},\mathrm{F})\) 细。反之，对拓扑 \(\tau(\mathrm{E},\mathrm{F})\) （引理 1）以及对拓扑 \(\sigma(\mathrm{E},\mathrm{F})\) （II, 第 43 页, 命题 3），F 都是 E 的对偶，因而对介于 \(\tau(\mathrm{E},\mathrm{F})\) 与 \(\sigma(\mathrm{E},\mathrm{F})\) 之间的每个拓扑，F 也是 E 的对偶。

推论. --- 设 p 是 E 上的一个半范数。下列条件等价：

\begin{enumerate}
\def\labelenumi{(\roman{enumi})}
\item
  \(p\) 对拓扑 \(\tau(\mathrm{E},\mathrm{F})\) 连续；
\item
  每个线性型 \(f\) （\(\mathbf{E}\) 上的）若满足 \(|f| \leqslant p\) ，都来自 \(\mathbf{F}\) 的一个元素。
\item[]
  \((i) \Rightarrow (ii):\) 若 \(p\) 对 \(\tau(\mathrm{E},\mathrm{F})\) 连续，则每个线性型 \(f\) （\(\mathrm{E}\) 上的）若满足 \(|f|\leqslant p\) ，就对 \(\tau(\mathrm{E},\mathrm{F})\) 连续，因而由引理 1 来自 \(\mathrm{F}\) 的一个元素。
\item[]
  \((ii) \Rightarrow (i):\) 设 T 是 E 上由半范数 p.~定义的拓扑。若条件 (ii) 满足，则 E 上对 T 连续的线性型都属于 F。由引理 2，T 比 \(\tau(\mathrm{E}, \mathrm{F})\) 粗，因而 p 对 \(\tau(\mathrm{E}, \mathrm{F})\) 连续。
\end{enumerate}

注记 2. --- * 设 K 是 F 的一个对弱拓扑 \(\sigma(F, E)\) 紧的凸子集，\(\mu\) 是 K 上的一个正测度。令

\[
p (x) = \int_ {\mathbf {K}} | \mathrm{B} (x, y) | d \mu (y)
\]

（对一切 \(x \in E\) ）。立即可见 p 是一个半范数。此外，对每个 \(x \in E\) ，关系式 \(\langle |B(x, y)| \leqslant 1 \text{ 对一切 } y \in K \rangle\) 蕴含 \(p(x) \leqslant \mu(K)\) 。这就证明 E 上的半范数 p 对 Mackey 拓扑 \(\tau(E, F)\) 连续。*

例. --- 设 G 是一个局部凸空间，\(G'\) 是它的对偶。在 \(G'\) 上，弱拓扑 \(\sigma(G', G)\) 与紧凸收敛拓扑（III, 第 14 页）都与 \(G'\) 和 G 之间的对偶相容。一般说来，\(G'\) 上的强拓扑与紧收敛拓扑都不与 \(G'\) 和 G 之间的对偶相容。然而回想，当 G 是 Hausdorff 的且拟完备的时候，\(G'\) 上的紧收敛拓扑与紧凸收敛拓扑一致（III, 第 8 页），因而与 \(G'\) 和 G 之间的对偶相容。

定义 3. --- 设 E 与 F 是两个处于对偶的向量空间，T 是 F 的对 \(\sigma(\mathrm{F}, \mathrm{E})\) 有界的子集所成的族。那么 F 上的 T-拓扑记作 \(\beta(\mathrm{E}, \mathrm{F})\) 。

类似地，我们在 F 上定义拓扑 \(\beta(\mathrm{F},\mathrm{E})\) 。容易看出，拓扑 \(\beta(\mathrm{E},\mathrm{F})\) 与 \(\beta(\mathrm{E},\mathrm{F}/\mathrm{E}^{\circ})\) 恒同，因而可以化归为 E 与 F 之间的对偶在 F 中是分离的情形。

注记. --- 3) 用 \(E_{\sigma}\) 表示赋以拓扑 \(\sigma(\mathrm{E}, \mathrm{F})\) 的空间 E。\(E_{\sigma}\) 中的桶（III, 第 24 页）是 E 的对 \(\sigma(\mathrm{E}, \mathrm{F})\) 为凸的、均衡的、闭的且吸收的子集。它们正是 F 的对 \(\sigma(\mathrm{F}, \mathrm{E})\) 为凸的、均衡的且有界的子集的极集。因此，\(E_{\sigma}\) 中所有桶的族是 E 中拓扑 \(\beta(\mathrm{E}, \mathrm{F})\) 的 0 的一个基本邻域系。换言之，E 上的一个半范数对 \(\beta(\mathrm{E}, \mathrm{F})\) 连续，当且仅当它对 \(\sigma(\mathrm{E}, \mathrm{F})\) 下半连续（cf.~III, 第 24 页, 命题 1）。

\begin{enumerate}
\def\labelenumi{\arabic{enumi})}
\setcounter{enumi}{3}
\item
  设 \(\mathcal{T}\) 是 E 上与 E 和 F 之间的对偶相容的一个拓扑。由 IV, 第 1 页, 命题 1, (ii)，F 上的拓扑 \(\beta(\mathrm{F},\mathrm{E})\) 正是 F 上的强拓扑（当把 F 与 E 的对偶（赋以拓扑 \(\mathcal{T}\) ）认同时）。
\item
  拓扑 \(\beta(\mathrm{E},\mathrm{F})\) （\(\mathrm{E}\) 上的）比 \(\tau(\mathrm{E},\mathrm{F})\) 细。一般说来，它与 \(\mathrm{E}\) 和 \(\mathrm{F}\) 之间的对偶不相容（但 cf.~§ 2）。特别地，\(\mathrm{E}\) 的一个对 \(\sigma(\mathrm{E},\mathrm{F})\) 有界的子集不一定对 \(\beta(\mathrm{E},\mathrm{F})\) 有界。
\end{enumerate}
